\pdfoutput=1
\documentclass[11pt]{amsart}
\usepackage{amsmath,amssymb,amsthm}
\usepackage{mathtools}
\usepackage{mathrsfs}
\usepackage[margin=1.2in]{geometry}
\usepackage{booktabs}
\usepackage{array}
\usepackage{flafter}
\usepackage{xcolor}
\usepackage[colorlinks=true,linkcolor=blue,citecolor=blue,urlcolor=blue]{hyperref}
\hypersetup{pdfauthor={Bernd Johannes Wuebben},
  pdftitle={One-bubble neighborhoods of Seiberg-Witten strata in SO(3)-monopole moduli spaces},
  pdfsubject={A family gluing and neighborhood theorem at the first Uhlenbeck level of the SO(3)-monopole moduli space},
  pdfkeywords={SO(3) monopoles, Uhlenbeck compactification, gluing,
  Seiberg-Witten moduli spaces, Donaldson invariants, Witten conjecture}}

\setlength{\emergencystretch}{2em}

\newtheorem{theorem}{Theorem}[section]
\newtheorem{proposition}[theorem]{Proposition}
\newtheorem{lemma}[theorem]{Lemma}
\newtheorem{corollary}[theorem]{Corollary}
\theoremstyle{definition}
\newtheorem{definition}[theorem]{Definition}
\newtheorem{convention}[theorem]{Convention}
\theoremstyle{remark}
\newtheorem{remark}[theorem]{Remark}

\theoremstyle{plain}
\newtheorem{thmletter}{Theorem}
\renewcommand{\thethmletter}{\Alph{thmletter}}
\newtheorem{corletter}[thmletter]{Corollary}

\newcommand{\RR}{\mathbb{R}}
\newcommand{\CC}{\mathbb{C}}
\newcommand{\cA}{\mathcal{A}}
\newcommand{\cB}{\mathcal{B}}
\newcommand{\cC}{\mathcal{C}}
\newcommand{\cD}{\mathcal{D}}
\newcommand{\cF}{\mathcal{F}}
\newcommand{\cG}{\mathcal{G}}
\newcommand{\cK}{\mathcal{K}}
\newcommand{\cM}{\mathcal{M}}
\newcommand{\cN}{\mathcal{N}}
\newcommand{\cU}{\mathcal{U}}
\newcommand{\cV}{\mathcal{V}}
\newcommand{\cX}{\mathcal{X}}
\newcommand{\cZ}{\mathcal{Z}}
\newcommand{\ad}{\operatorname{ad}}
\newcommand{\coker}{\operatorname{coker}}
\newcommand{\diag}{\operatorname{diag}}
\newcommand{\Fr}{\operatorname{Fr}}
\newcommand{\Gl}{\operatorname{Gl}}
\newcommand{\id}{\operatorname{id}}
\newcommand{\im}{\operatorname{im}}
\newcommand{\ind}{\operatorname{ind}}
\newcommand{\rank}{\operatorname{rank}}
\newcommand{\SO}{\operatorname{SO}}
\newcommand{\SU}{\operatorname{SU}}
\newcommand{\SW}{\operatorname{SW}}

\title[One-bubble neighborhoods for $\SO(3)$ monopoles]
{One-bubble neighborhoods of Seiberg--Witten strata\\
in $\SO(3)$-monopole moduli spaces}
\author{Bernd Johannes Wuebben}

\date{30 August 2026}
\subjclass[2020]{57K41 (primary); 57R57, 58J05, 58J37, 53C07 (secondary)}
\keywords{$SO(3)$ monopoles, Uhlenbeck compactification, gluing,
Seiberg--Witten moduli spaces, Donaldson invariants, Witten conjecture}

\begin{document}

\begin{abstract}
Feehan and Leness constructed local gluing and obstruction maps near
Seiberg--Witten strata in the Uhlenbeck compactification of the
$SO(3)$-monopole moduli space.  Their level-one link calculation also requires
continuity, embedding, surjectivity onto a Uhlenbeck-open neighborhood, and
orientation properties
assigned to a companion gluing theorem and retained as a hypothesis.  We
prove these properties at the one-bubble level over any compact regular
Seiberg--Witten moduli space consisting of solutions with nonzero spinor.

A global equivariant stabilization of the normal deformation operators
handles the family problem without a fixed spectral cutoff.  A scale-uniform
augmented parametrix in mixed critical-exponent norms gives
$O(\lambda^{1/3})$ bounds for the projected splicing error and nonlinear
correction, uniformly over the family.  Projection onto fixed-model Jacobi
fields recovers the background and gluing parameters; a two-sided spectral
estimate on the conformal neck gives the no-neck theorem needed for reverse
gluing.  At zero scale the instanton obstruction line collapses to the vertex
of a cone family, and equivariant determinant-line excision identifies the
actual and virtual links as oriented cycles.

Consequently, this proves the neighborhood-theorem input to the
Feehan--Leness level-one pairing over regular strata of arbitrary dimension.
Substitution into their published cohomological calculation yields its
Jacobi-polynomial formula; the zero-dimensional formula and the associated
low-degree Donaldson--Seiberg--Witten comparison follow as corollaries.  The
result is confined to one bubble and does not treat collision strata,
zero-spinor solutions, or the full Witten conjecture.
\end{abstract}

\maketitle
\enlargethispage{6pt}

\section{Introduction}\label{sec:introduction}

\subsection{The neighborhood problem}
Let $X$ be a smooth, closed, connected, oriented four-manifold, and consider
the Uhlenbeck compactification of a moduli space of $SO(3)$ monopoles.  Its
lower levels contain products of Seiberg--Witten moduli spaces with symmetric
products of $X$.  The $SO(3)$-monopole cobordism program compares the link of
the top anti-self-dual stratum with the links of these reducible strata.  A
link calculation has two distinct analytic inputs.  One must first construct
solutions of a finite-dimensional extension of the monopole equations from
splicing data.  One must then prove that the resulting zero set is an actual
neighborhood of the ideal stratum in the Uhlenbeck topology.

The first input was established in the local gluing theorem of Feehan and
Leness~\cite{FeehanLenessIII}.  That theorem constructs the gluing-data
bundle, a smooth solution map for the extended equations, and a finite-rank
obstruction section.  It gives the inclusion of the obstruction zero set in
the monopole moduli space.  Section~1.5.1 of the same paper separates from
this construction the remaining continuity and embedding statements and the
claim of surjectivity onto a Uhlenbeck-open neighborhood; it assigns this
package to a companion article.  In the one-bubble case,
\cite[Theorem~3.8]{FeehanLenessLevelOne} records the compactified embedding,
the obstruction sections, and the inclusion of their zero set in the monopole
moduli space, but it does not assert that this zero set maps onto an open
neighborhood of the ideal stratum.  That stronger conclusion is Property~(4)
of Hypothesis~7.8.1 in~\cite{FeehanLenessCobordism} and is retained there as
part of the local gluing hypothesis.

This paper proves the missing neighborhood statement for one charge-one
bubble over a compact regular Seiberg--Witten moduli space of arbitrary
dimension.  At this level there are no collision strata and no overlaps
among partitions of a symmetric product.  The family problem cannot be
reduced to the zero-dimensional case by choosing one spectral cutoff:
spectral flow can change the rank of the corresponding small-eigenvalue
spaces.  We instead stabilize the normal family globally by a smooth
finite-rank equivariant bundle.  Local cutoff reductions enter only through
a common enlargement, so neither the gluing chart nor its orientation
depends on a constant-rank spectral window.

\subsection{Results}
Let $b_1(X)=0$ and let $b_2^+(X)\geq3$ be odd.  Fix generic perturbation data,
a $\operatorname{spin}^u$ structure $\mathfrak t'$, and a
$\operatorname{spin}^c$ structure $\mathfrak s$ for which the level-one
reducible stratum is
\begin{equation}\label{eq:intro-stratum}
 M_{\mathfrak s}\times X
 \subset \overline{\cM}_{\mathfrak t'}.
\end{equation}
Let $\mathfrak t$ denote the lower-level $\operatorname{spin}^u$ structure
containing $M_{\mathfrak s}$.  In the standard level-one notation,
\begin{equation}\label{eq:lower-ambient-structures}
 \mathfrak t=\mathfrak s\oplus\mathfrak s\otimes L,
 \qquad c_1(\mathfrak t)=c_1(\mathfrak t'),
 \qquad p_1(\mathfrak t)=p_1(\mathfrak t')+4.
\end{equation}
Thus $\mathfrak t'$ is the ambient structure obtained after splicing one unit
of instanton charge into a configuration for $\mathfrak t$.
Assume that $M_{\mathfrak s}$ is a compact regular smooth manifold, every
represented Seiberg--Witten pair has nonzero spinor, and
\begin{equation}\label{eq:intro-sw-dimension}
 \dim M_{\mathfrak s}=d_{\mathfrak s}=2d\geq0.
\end{equation}
The level-one construction has a compactified
virtual gluing space
\begin{equation}\label{eq:intro-virtual}
 \overline{\cM}^{\mathrm{vir}}_{\mathfrak t',\mathfrak s}
 (\varepsilon,\delta),
\end{equation}
an obstruction pseudo-bundle $\overline\Upsilon$, and a section $\chi$.
Their definitions are recalled in Section~\ref{sec:gluing-data}.  The
instanton summand of $\overline\Upsilon$ has complex rank one at positive
scale and collapses at the zero-scale cone point.

Writing $N(\varepsilon)$ for the background virtual-normal disk, the
compactified gluing space has the three strata
\begin{equation}\label{eq:intro-three-strata}
 \cG^+\quad\sqcup\quad
 \underbrace{(N(\varepsilon)-M_{\mathfrak s})\times X}_{\text{middle
 zero-scale stratum}}\quad\sqcup\quad
 \underbrace{M_{\mathfrak s}\times X}_{\text{lowest stratum}}.
\end{equation}
At positive scale the obstruction section has background and instanton
components $\chi=(\chi_s,\chi_i)$.  At zero scale the instanton fiber and
$\chi_i$ collapse, while $\chi_s$ restricts to the background obstruction
section $b:N(\varepsilon)\to\Xi$.  Thus $b$ is transverse on the middle
stratum and vanishes on the lowest stratum.

\begin{thmletter}[level-one neighborhood; $=$
Theorem~\ref{thm:neighborhood}]\label{thm:intro-neighborhood}
For sufficiently small $\varepsilon,\delta>0$, there is an
$S^1$-equivariant map
\[
 \gamma:\overline{\cM}^{\mathrm{vir}}_{\mathfrak t',\mathfrak s}
 (\varepsilon,\delta)\longrightarrow\overline\cC_{\mathfrak t'}
\]
with the following properties.
\begin{enumerate}
\item The map is a topological embedding, is smooth on the positive-scale
stratum, and is joined to the splicing map by an $S^1$-equivariant isotopy
through compactified embeddings.  On the zero-scale face it is the canonical
background-and-ideal-point map.
\item There is an Uhlenbeck-open neighborhood $\cU$ of
$M_{\mathfrak s}\times X$ for which the restriction
\[
 \gamma:\chi^{-1}(0)\cap\gamma^{-1}(\cU)\longrightarrow
 \cU\cap\overline\cM_{\mathfrak t'}^{*,0}
\]
is a homeomorphism and is a diffeomorphism on the positive-scale stratum.
\item The total obstruction section is transverse on the positive-scale
stratum, and its background component is transverse on the middle
zero-scale stratum.  Its norm is continuous along the compactified section.
No transversality is asserted on the lowest stratum.
\item For a homology orientation of $X$, determinant-line excision gives the
orientation of the transverse zero locus.  This orientation descends through
the frame and residual-circle quotients and induces the standard boundary
orientation used in the level-one link pairing.
\end{enumerate}
\end{thmletter}

The last qualification in (3) is forced by the geometry.  On the lowest
stratum the surviving background obstruction section is identically zero.
If its rank is positive, its restriction cannot be transverse.  This is the
statement used in the relative-Euler-class calculation of
\cite[Section~5]{FeehanLenessLevelOne}; it also agrees with the qualification
following Theorem~3.8 there.  Under the fine decomposition
\eqref{eq:intro-three-strata}, the all-strata transversality formulation in
\cite[Hypothesis~7.8.1(2)]{FeehanLenessCobordism} would also include the
lowest stratum and is false there when the background obstruction has
positive rank.

The analytic estimate behind Theorem~\ref{thm:intro-neighborhood} may be
stated independently.  Write $q_\lambda$ for the standard one-bubble
pregluing and $J_{s,\lambda}$ for the matched embedding of the background
obstruction space.  If $\mathfrak S$ is the monopole map and
$\mathfrak S(q)=\iota_{s,q}b(q)$ on the background virtual neighborhood,
then
\begin{equation}\label{eq:intro-residual}
 \left\|\mathfrak S(q_\lambda)-J_{s,\lambda}b(q)\right\|_{L^{\sharp,2;2}}
 \leq C\lambda^{1/3}.
\end{equation}
There is a uniformly bounded inverse for the augmented linearization on a
model-coordinate complement of all gluing parameters.  The contraction
therefore gives a unique correction and obstruction coefficients satisfying
\begin{equation}\label{eq:intro-correction}
 \|u_\lambda\|_{L^2_1}
 +\|\eta_{s,\lambda}\|+\|\eta_{i,\lambda}\|
 \leq C\lambda^{1/3}.
\end{equation}

Theorem~\ref{thm:intro-neighborhood} is the principal result.  The next two
statements are consequences obtained by inserting it into published
cohomological calculations.  The neighborhood theorem supplies every
invocation of the deferred gluing result in Sections~3--6 of
\cite{FeehanLenessLevelOne}.  We retain the notation of that paper for the
pairing.  Let $\mathfrak t'$ be a
$\operatorname{spin}^u$ structure with $w_2(\mathfrak t')$ good in the sense
of~\cite[Definition~2.3]{FeehanLenessLevelOne}, and suppose
\begin{equation}\label{eq:intro-level-relation}
 \bigl(c_1(\mathfrak t')-c_1(\mathfrak s)\bigr)^2
 =p_1(\mathfrak t')+4.
\end{equation}
Let $h\in H_2(X;\RR)$ and let $x\in H_0(X;\mathbb Z)$ be the positive
generator.  For a monomial $z$ in $h$ and $x$, let $\overline V(z)$ denote
the closure of a transverse geometric representative of the Donaldson class
$\mu_p(z)$; let $\overline W$ represent the circle class $\mu_c$; and let
$\overline L_{\mathfrak t',\mathfrak s}$ be the standard oriented link of
$M_{\mathfrak s}\times X$ in
$\overline\cM_{\mathfrak t'}^{*,0}/S^1$.  The symbol $\#$ below denotes the
signed count of a transverse zero-dimensional intersection.  Let
$\mu_{\mathfrak s}\in H^2(M_{\mathfrak s};\mathbb Z)$ be the first Chern
class of the Seiberg--Witten base-point fibration.  With the orientation
determined by the chosen homology orientation, set
\begin{equation}\label{eq:sw-definition}
 \SW_X(\mathfrak s)
 =\left\langle\mu_{\mathfrak s}^{d},[M_{\mathfrak s}]\right\rangle.
\end{equation}
Also write
\begin{equation}\label{eq:asd-dimension}
 d_a(\mathfrak t')=-2p_1(\mathfrak t')
 -\tfrac32\bigl(\chi(X)+\sigma(X)\bigr),
\end{equation}
and set
\begin{equation}\label{eq:intro-pairing-notation}
 \beta=c_1(\mathfrak s)-c_1(\mathfrak t'),\qquad
 B_h=\langle\beta,h\rangle.
\end{equation}
For the Jacobi polynomials we use the normalization
\begin{equation}\label{eq:intro-jacobi-normalization}
 P_d^{A,B}(\zeta)=2^{-d}\sum_{v=0}^d
 \binom{d+A}{v}\binom{d+B}{d-v}
 (\zeta-1)^{d-v}(\zeta+1)^v.
\end{equation}

\begin{corletter}[completion of the family level-one pairing; $=$
Corollary~\ref{cor:pairing}]\label{cor:intro-pairing}
Let $\delta,m,\eta$ be nonnegative integers satisfying
\[
 0\leq m\leq\lfloor\delta/2\rfloor,
 \qquad
 2(\delta+\eta)=\dim(\cM_{\mathfrak t'}/S^1)-1.
\]
Put $n=\delta-2m$ and
\begin{align*}
 A_J&=\eta-d+1,\qquad
 B_J=\frac12\bigl(2\delta-d_a(\mathfrak t')-d_{\mathfrak s}\bigr)
      -\frac14\bigl(\chi(X)+\sigma(X)\bigr),\\
 J_d&=P_d^{A_J,B_J}(0),\qquad
 J_d'=P_d^{A_J-1,B_J+1}(0).
\end{align*}
With the standard level-one link orientation,
\begin{align*}
&\#\bigl(
\overline V(h^{\delta-2m}x^m)\cap\overline W^\eta
\cap\overline L_{\mathfrak t',\mathfrak s}\bigr)\\
&\quad=(-1)^{m+1+d}2^{d-\delta}
 \left\langle\mu_{\mathfrak s}^{d},[M_{\mathfrak s}]\right\rangle\\
&\qquad\quad\times
 \left[J_d\left(a_0B_h^n+a_1B_h^{n-2}Q_X(h,h)\right)
 +2nJ_d'B_h^{n-1}\langle c_1(\mathfrak t'),h\rangle\right],
\end{align*}
where terms with negative powers of $B_h$ are omitted and
\begin{align*}
a_0&=3\beta^2+c_1^2(X)+2\beta\cdot c_1(\mathfrak t')
     +4n-4m,\qquad
a_1=4\binom n2.
\end{align*}
\end{corletter}

This is the polynomial form of~\cite[Theorem~6.1]{FeehanLenessLevelOne}; it
remains meaningful when $J_d=0$.  The new neighborhood theorem supplies the
deferred analytic hypothesis of that calculation.  Its universal-class
pullbacks, equivariant Thom and Segre reduction, charge-one fiber integrals,
and Jacobi identity remain published inputs.  When $d=0$, one has
$J_0=J_0'=1$ and the theorem reduces exactly to the zero-dimensional formula
proved in the earlier version of this manuscript.

For completeness, write
\[
 c(X)=-\frac14\bigl(7\chi(X)+11\sigma(X)\bigr).
\]
Combining Corollary~\ref{cor:intro-pairing} with the published top-level
calculation, blow-up formulas, and superconformal-simple-type vanishing gives
the following precise consequence.

\begin{corletter}[degree consequence; $=$
Corollary~\ref{cor:degree}]\label{cor:intro-degree}
Suppose in addition that $X$ is abundant, effective, and of
Seiberg--Witten simple type, and choose the characteristic data as in
\cite[Theorem~1.1]{FeehanLenessLevelOne}.  Then
\[
 \cD_X^w(h)\equiv0\equiv\SW_X^w(h)\pmod{h^{c(X)-2}}
\]
and
\[
 \cD_X^w(h)\equiv
 2^{2-c(X)}e^{Q_X(h,h)/2}\SW_X^w(h)
 \pmod{h^{c(X)+2}}.
\]
\end{corletter}

The corollary retains every hypothesis just listed.  It does not prove the
full Witten conjecture, remove effectiveness, or treat a level
$\ell\geq2$ contribution.

\subsection{Method}
There are four analytic points.  First, the compact family of background
normal operators is stabilized by a global complex equivariant bundle
$\Xi$.  The resulting augmented operator contains both $\Xi$ and the complex
one-dimensional cokernel of the charge-one Dirac operator.  A logarithmic
square partition gives a cutoff error
$C(\log N)^{-1/2}$.  The physical bubble geometry is compactified to a family
of metrics on $S^4$; using the induced operator eliminates a principal-symbol
comparison which is not bounded in the critical target.  All remaining
coefficient errors are $O_N(\lambda^{1/3})$.  The inverse and its first two
normalized parameter derivatives are uniformly bounded on transported graph
domains, which gives a smooth correction over the family.

Second, the complement is not the physical $L^2$-orthogonal complement of
the bubble Jacobi fields.  Their physical $L^2$ norms collapse with scale.
Instead, a glued variation is localized, transported to the fixed background
and $S^4$ models, and projected there.  Synthesis followed by this retraction
is $I+o(1)$.  It gives a uniform projection onto the background normal,
translation, dilation, and relative-frame modes and a uniformly coercive
complement.

Third, parameter recovery requires a slice estimate uniform both in the
family and in the bubbling scale.  We prove a critical Green estimate for
the coupled pair gauge Laplacian.  Exterior-based gauge normalization
retains the constant adjoint mode at bubble infinity as the relative frame;
normalizing at a finite point on the bubble would erase precisely that
gluing parameter.

Fourth, continuity and compactness do not imply neighborhood surjectivity.
For a genuine monopole sequence converging to the ideal stratum, a
charge-one rescaling and a fixed-annulus estimate reduce the neck to a
first-order equation on a conformal cylinder.  The nonzero spectrum of the
odd-signature and Dirac tangential operators is separated uniformly from
zero.  A two-sided Duhamel estimate then gives a length-independent
$\ell^1$ bound, including the $r^{3/2}$ outer spinor-cutoff term.  Exact
modulation places the resulting difference in the complement used for the
construction.  The augmented inverse and contraction uniqueness identify
the original monopole with the gluing solution.

The determinant calculation follows the same decomposition.  Excision splits
the background and bubble determinant lines.  The complex background normal
and obstruction factors cancel in the finite-dimensional reduction, as do
the tangent- and adjoint-frame quotient factors.  The residual-circle and
outward-radial conventions give
\begin{equation}\label{eq:intro-orientation}
 \lambda_{\cM}=-v_{S^1}\wedge\vec r\wedge\widetilde\lambda_L,
\end{equation}
the convention in~\cite[Equation~(3.81)]{FeehanLenessLevelOne}.

\subsection{Relation to prior work}
The local construction quoted here is the one-bubble specialization of
\cite{FeehanLenessIII}.  The critical norms and uniform slice estimates are
from~\cite{FeehanCritical}; the Uhlenbeck gauge theorem and compactness
principles are used in their standard form~\cite{Uhlenbeck82,DK90}.  The
induced-metric construction follows the conformal comparison in
\cite[Section~8]{FeehanLenessIII}.  The parameter-recovery and reverse-gluing
arguments are adapted to the coupled monopole equation rather than imported
from ASD gluing: the bubble spinor, mixed critical target, and background
obstruction must be controlled simultaneously.

Theorem~\ref{thm:intro-neighborhood} concerns the first lower Uhlenbeck level
only.  The topological splicing and overlap problem at all levels is treated
in~\cite{FeehanLenessCobordism}, while the analytic neighborhood
identification is assumed there.  At level one the overlap problem is absent,
which is what makes it possible to separate and prove the neighborhood
theorem without solving the collision problem.  The varying-rank issue is
handled here by global family stabilization and stable comparison with the
local spectral reductions, not by assuming that a spectral window has
constant rank.

\section{Level-one gluing data}\label{sec:gluing-data}

\subsection{The monopole map and its deformation operator}
Let $\mathfrak t'$ be represented by a Clifford module
$V=V^+\oplus V^-$ with Clifford multiplication $\rho$ and adjoint bundle
$\mathfrak g_V$.  A configuration is a pair $(A,\Phi)$ consisting of a
compatible connection and a section of $V^+$.  We use the perturbed monopole
map
\begin{equation}\label{eq:monopole-map}
 \mathfrak S(A,\Phi)=
 \left(
 F_A^+-\tau\rho^{-1}(\Phi\otimes\Phi^*)_{00},
 (D_A+\rho(\vartheta))\Phi
 \right).
\end{equation}
The perturbations are fixed throughout and belong to a compact generic
family.  For the determinant-one gauge group, our infinitesimal-action
convention is
\[
 d^0_{A,\Phi}\zeta=(d_A\zeta,-\zeta\Phi).
\]
Thus
\[
 d^{0,*}_{A,\Phi}(a,\phi)=d_A^*a-(\,\cdot\,\Phi)^*\phi,
\]
and the derivative of~\eqref{eq:monopole-map} is
\begin{align}\label{eq:d1}
d^1_{A,\Phi}(a,\phi)=\bigl(&d_A^+a-
\tau\rho^{-1}(\Phi\otimes\phi^*+\phi\otimes\Phi^*)_{00},\notag\\
&(D_A+\rho(\vartheta))\phi+\rho(a)\Phi\bigr).
\end{align}
The gauge-fixed operator is
\begin{equation}\label{eq:gauge-fixed}
 \mathscr D_{A,\Phi}=d^{0,*}_{A,\Phi}\oplus d^1_{A,\Phi}.
\end{equation}
At an approximate solution $d^1d^0$ need not vanish.  All parametrices below
include the gauge row and do not use a deformation complex at the preglued
pair.

Equation~(2.3) of~\cite{FeehanLenessIII} prints a factor $1/2$ in the
spinor-bilinear term of $d^1$, whereas the literal derivative of the displayed
monopole map, Equation~(9.3) there, and
\cite[Equation~(2.36)]{FeehanLenessI} have the normalization in
\eqref{eq:d1}.  We use~\eqref{eq:d1}.  Interpolating the discrepant
zeroth-order coefficient through $[1/2,1]$ does not change a determinant-line
orientation; the characteristic-class normalization is checked independently
in Section~\ref{sec:pairing}.

\subsection{The background finite-dimensional reduction}
The lower-level structure in~\eqref{eq:lower-ambient-structures} has the
reducible splitting
\[
 \mathfrak t=\mathfrak s\oplus\mathfrak s\otimes L.
\]
Along the reducible embedding of $M_{\mathfrak s}$, the gauge-fixed family
splits into tangential and normal summands
\begin{equation}\label{eq:family-splitting}
 \mathscr D=\mathscr D^t\oplus\mathscr D^n.
\end{equation}
Regularity of the Seiberg--Witten equations gives
\begin{equation}\label{eq:tangential-family}
 \ker\mathscr D^t=TM_{\mathfrak s},\qquad
 \coker\mathscr D^t=0.
\end{equation}
The normal family is complex linear, and the residual circle acts on its
domain and target through one complex character.

We use the extrinsic finite-dimensional reduction of
\cite[Theorem~3.19]{FeehanLenessTopLevel}.  Compactness of $M_{\mathfrak s}$
gives a finite-rank complex bundle $\Xi$ and an equivariant target embedding
$\iota_\Xi$ such that
\begin{equation}\label{eq:global-augmentation}
 \widehat{\mathscr D}^{n}_p
 =\mathscr D^n_p\oplus(-\iota_{\Xi,p}):
 \operatorname{Dom}\mathscr D^n_p\oplus\Xi_p
 \longrightarrow\operatorname{Tar}\mathscr D^n_p
\end{equation}
is surjective for every $p\in M_{\mathfrak s}$.  Indeed, each $p$ admits a
finite-dimensional complex cokernel-spanning space.  Surjectivity persists
on a neighborhood, a finite subcover supplies a common enlargement, and an
equivariant tubular retraction extends it to a small normal neighborhood.
No spectral projection is used in this construction.

The augmented kernels form a smooth complex bundle
\begin{equation}\label{eq:normal-bundle-definition}
 N_{\mathfrak t,\mathfrak s}
 =\ker\widehat{\mathscr D}^{n}\longrightarrow M_{\mathfrak s}.
\end{equation}
Thus the full augmented kernel is
$TM_{\mathfrak s}\oplus N_{\mathfrak t,\mathfrak s}$ and
\begin{equation}\label{eq:index-family}
 \operatorname{Ind}_{\CC}(\mathscr D^n)
 =[N_{\mathfrak t,\mathfrak s}]-[\Xi]
 \quad\text{in }K^0(M_{\mathfrak s}).
\end{equation}
The projected convention of~\cite{FeehanLenessIII} is equivalent: if
$\Pi_{\Xi^perp}$ is orthogonal projection, then
\begin{equation}\label{eq:projected-augmented-kernel}
 u\longmapsto\bigl(u,\iota_\Xi^{-1}\mathscr D^nu\bigr)
\end{equation}
identifies $\ker(\Pi_{\Xi^\perp}\mathscr D^n)$ with
$\ker\widehat{\mathscr D}^{n}$.

After choosing a sufficiently small lifted disk in
$N_{\mathfrak t,\mathfrak s}$, the implicit-function theorem gives a smooth
embedding into the thickened configuration quotient and a section
\begin{equation}\label{eq:background-obstruction}
 b:N_{\mathfrak t,\mathfrak s}(\varepsilon)\longrightarrow\Xi,
 \qquad
 \mathfrak S(q)=\iota_{s,q}b(q).
\end{equation}
Its zero set is $M_{\mathfrak s}$.  The term \emph{virtual-normal disk}
will always mean this disk bundle; it never means an eigenspace selected by
a spectral cutoff.  Since every pair on the zero section has nonzero spinor,
compactness and openness allow the disk to be chosen so that every pair in
it has trivial determinant-one infinitesimal stabilizer.  Put
\[
 r_N=\rank_{\CC}N_{\mathfrak t,\mathfrak s},\qquad
 r_\Xi=\rank_{\CC}\Xi.
\]
The background parameter bundle therefore has real dimension
$d_{\mathfrak s}+2r_N$, while its obstruction variables have real dimension
$2r_\Xi$.

\subsection{The charge-one factor and compactification}
Let $M_1^{s,\natural}(S^4,\delta)$ be the centered framed charge-one ASD
moduli space with scale in $(0,\delta]$.  There are natural identifications
\begin{equation}\label{eq:charge-one-cone}
 M_1^{s,\natural}(S^4,\delta)
 \cong(0,\delta]\times\SO(3),
 \qquad
 \overline M_1^{s,\natural}(S^4,\delta)
 \cong c(\SO(3)).
\end{equation}
If $I$ is the standard centered instanton, then
\begin{equation}\label{eq:bubble-cokernel}
 \ker D_I=0,
 \qquad K_I:=\coker D_I=\ker D_I^*,
 \qquad \dim_{\CC}K_I=1.
\end{equation}
The framed ASD operator is surjective.  Its eight-dimensional kernel consists
of four translations, one dilation, and three framed rotations.

The pre-splicing factor is
\begin{align}\label{eq:gluing-quotient}
\widetilde\Gl_{\mathfrak t}(\delta)
=\bigl(&\Fr(\mathfrak g_V)\times_X\Fr(T^*X)
\times\Fr(\mathfrak g_V|_s)\notag\\
&\times\widetilde M_1^\natural(S^4,\delta)\bigr)
/(\SO(3)\times\SO(4)).
\end{align}
After dividing by the charge-one gauge group, one obtains
$\Gl_{\mathfrak t}(\delta)$.  Replacing the instanton factor by the cone in
\eqref{eq:charge-one-cone} gives $\overline\Gl_{\mathfrak t}(\delta)$ and
\begin{equation}\label{eq:virtual-space}
 \overline\cM^{\mathrm{vir}}_{\mathfrak t',\mathfrak s}
 =\widetilde N_{\mathfrak t,\mathfrak s}(\varepsilon)
 \times_{\cG_{\mathfrak s}}\overline\Gl_{\mathfrak t}(\delta).
\end{equation}
At the cone point the relative instanton frame is forgotten.  Hence
\eqref{eq:virtual-space} has three strata:
\begin{align}\label{eq:three-strata}
\overline\cM^{\mathrm{vir}}_{\mathfrak t',\mathfrak s}
={}&\cG^+\notag\\
&\sqcup\bigl((N_{\mathfrak t,\mathfrak s}(\varepsilon)
-M_{\mathfrak s})\times X\bigr)
\sqcup(M_{\mathfrak s}\times X),
\end{align}
where $\cG^+$ denotes the positive-scale stratum.

Before quotienting, the splicing map is equivariant for the diagonal
$\cG_{\mathfrak s}$ and charge-one gauge actions and for the diagonal
$\SO(4)$ tangent-frame and $\SO(3)$ adjoint-frame actions.  These auxiliary
frame actions are divided out; the relative $\SO(3)$ frame in
\eqref{eq:charge-one-cone} remains.

The residual circle has the two equivalent presentations of
\cite[Lemma~3.6]{FeehanLenessLevelOne}.  It may act by the monopole circle on
the normal factor and trivially on the compactified gluing factor, or
diagonally with weight two on both.  We use the effective circle in the link
quotient.  Passing from weight two to its effective parametrization does not
reverse the orbit orientation.

\subsection{The obstruction pseudo-bundle}
At positive scale the obstruction space is
\begin{equation}\label{eq:obstruction-fiber}
 \Upsilon_g=\Upsilon_g^s\oplus\Upsilon_g^i
 \cong\Xi_q\oplus K_I.
\end{equation}
The background summand extends as an ordinary bundle over
\eqref{eq:three-strata}.  The instanton summand is the associated line
\[
 (0,\delta]\times\SU(2)\times_{\{\pm1\}}\CC
 \longrightarrow (0,\delta]\times\SO(3).
\]
Its compactification is the cone map
\begin{equation}\label{eq:obstruction-cone}
 c(\SU(2)\times_{\{\pm1\}}\CC)
 \longrightarrow c(\SO(3));
\end{equation}
the fiber collapses at the vertex.  This is a rank-dropping pseudo-bundle,
not a singular trivialization of a line bundle.

\subsection{Strong and critical-norm Sobolev completions}
For smooth dependence at a fixed positive scale, choose $k\geq4$ and use the
standard Fredholm map
\[
 \mathscr D_{q_\lambda}:L^2_k(\Lambda^1\otimes\mathfrak g_V\oplus V^+)
 \longrightarrow
 L^2_{k-1}(\mathfrak g_V\oplus\Lambda^+\otimes\mathfrak g_V\oplus V^-).
\]
Uniform estimates as $\lambda\downarrow0$ use
\begin{align}\label{eq:weak-spaces}
 \mathbb X_{q_\lambda}&=L^2_{1,A_\lambda},\notag\\
 \mathbb Z_{q_\lambda}&=L^2(\mathfrak g_V),\notag\\
 \mathbb Y_{q_\lambda}&=
 L^{\sharp,2}(\Lambda^+\otimes\mathfrak g_V)\oplus L^2(V^-).
\end{align}
Here $L^{\sharp,2}=L^\sharp\cap L^2$ with the critical norm of
\cite{FeehanCritical}.  We construct the inverse on smooth target data, prove
the estimate in~\eqref{eq:weak-spaces}, and extend by density.  We do not
assert that the full first-order operator maps every arbitrary $L^2_1$
section into the mixed critical target.

Choose matched finite-dimensional embeddings
\[
 J_{s,\lambda}:\Xi_q\to\mathbb Y_{q_\lambda},\qquad
 J_{i,\lambda}:K_I\to\mathbb Y_{q_\lambda}.
\]
The augmented linearization is
\begin{equation}\label{eq:augmented-operator}
\widehat{\mathscr D}_\lambda(u,\xi_s,\xi_i)
=\left(d^{0,*}_{q_\lambda}u,
d^1_{q_\lambda}u-J_{s,\lambda}\xi_s-J_{i,\lambda}\xi_i\right).
\end{equation}
The nonlinear augmented map is
\begin{align}\label{eq:augmented-map}
\widehat\cF_\lambda(u,\eta_s,\eta_i)
=\bigl(&d^{0,*}_{q_\lambda}u,
\mathfrak S(q_\lambda+u)\notag\\
&-J_{s,\lambda}(b(q)+\eta_s)-J_{i,\lambda}\eta_i\bigr).
\end{align}

The retained parameter directions have real dimension
$d_{\mathfrak s}+2r_N+8$: the background tangent and normal variables, four
translations, dilation, and three relative rotations.  The obstruction rank
is $2r_\Xi+2$.  Hence the expected dimension of the obstruction zero set is
\begin{equation}\label{eq:dimension-check}
 d_{\mathfrak s}+(2r_N+8)-(2r_\Xi+2)
 =d_{\mathfrak s}+2(r_N-r_\Xi)+6.
\end{equation}

\begin{theorem}[level-one neighborhood]\label{thm:neighborhood}
Let $b_1(X)=0$ and let $b_2^+(X)\geq3$ be odd.  Fix generic perturbation data
and a level-one pair $(\mathfrak t',\mathfrak s)$ as in
\eqref{eq:intro-stratum}--\eqref{eq:lower-ambient-structures}.  Suppose that
$M_{\mathfrak s}$ is compact and regular and that every represented
Seiberg--Witten pair has nonzero spinor.  For sufficiently small
$\varepsilon,\delta>0$, there is
an equivariant compactified embedding
\[
 \gamma:\overline{\cM}^{\mathrm{vir}}_{\mathfrak t',\mathfrak s}
 (\varepsilon,\delta)\longrightarrow\overline\cC_{\mathfrak t'}
\]
and a continuous compactified obstruction section
$\chi=(\chi_s,\chi_i)$ such that the following hold: the positive-scale maps
are smooth; $\gamma$ is equivariantly isotopic through embeddings to
splicing; $\chi^{-1}(0)$ maps homeomorphically onto an Uhlenbeck-open
neighborhood of $M_{\mathfrak s}\times X$ and diffeomorphically on the
positive-scale stratum; $\chi$ is transverse there and $\chi_s$ is transverse
on the middle zero-scale stratum; and the determinant-line orientation
induces the standard level-one boundary orientation.
\end{theorem}

The proof occupies Sections~\ref{sec:parametrix}--\ref{sec:orientations}.
The local pregluing construction and the one-sided inclusion
$\gamma(\chi^{-1}(0))\subset\overline\cM_{\mathfrak t'}^{*,0}$ are quoted
from~\cite{FeehanLenessIII}.  Everything needed to turn that inclusion into
the neighborhood homeomorphism is proved below.

\section{The uniform augmented parametrix}\label{sec:parametrix}

\subsection{Critical norms and the square partition}
For a bundle-valued section $f$ on $X$, put
\begin{align}\label{eq:critical-norms}
 \|f\|_{L^\sharp}
 &=\sup_{z\in X}\int_X\operatorname{dist}(z,y)^{-2}|f(y)|\,dV(y),\notag\\
 \|f\|_{L^{2\sharp}}
 &=\sup_{z\in X}
 \left(\int_X\operatorname{dist}(z,y)^{-2}|f(y)|^2\,dV(y)\right)^{1/2},\\
 \|f\|_{L^{\sharp,2}}&=\|f\|_{L^\sharp}+\|f\|_{L^2}.
 \notag
\end{align}
The critical multiplication and Sobolev estimates of
\cite[Lemmas~4.1 and~4.3]{FeehanCritical} give
\begin{equation}\label{eq:critical-products}
 L^2_1\hookrightarrow L^{2\sharp},\qquad
 L^{2\sharp}\cdot L^{2\sharp}\longrightarrow L^\sharp,
\end{equation}
uniformly for the compact family of background connections under
consideration.  In particular,
\begin{equation}\label{eq:critical-sobolev}
 \|v\|_{L^4}+\|v\|_{L^{2\sharp}}
 \leq C\|v\|_{L^2_{1,A}}.
\end{equation}

Fix $N>4$.  The logarithmic cutoff of
\cite[Lemma~5.4]{FeehanLenessIII}, pulled back by geodesic normal coordinates
at $x$, may be chosen to satisfy
\begin{align}\label{eq:log-cutoff}
 \beta&=0 &&\text{if }r\leq N^{-1}\sqrt\lambda,
 &\beta&=1 &&\text{if }r\geq\tfrac12\sqrt\lambda,\notag\\
 \|d\beta\|_{L^4}&\leq C(\log N)^{-3/4},
 &\|d\beta\|_{L^{2\sharp}}&\leq C(\log N)^{-1/2},
 &\|d\beta\|_{L^2}&\leq C(\log N)^{-1}\sqrt\lambda.
\end{align}
Define
\begin{equation}\label{eq:square-partition}
 \beta_X=\sin(\pi\beta/2),\qquad
 \beta_I=\cos(\pi\beta/2).
\end{equation}
Then $\beta_X^2+\beta_I^2=1$, and both derivatives obey the
bounds in~\eqref{eq:log-cutoff}.  These estimates are uniform in $x$: near
$x$ this follows from normal-coordinate comparison, and for a base point
outside a fixed larger chart the last, scale-decaying bound in
\eqref{eq:log-cutoff} controls the supremum in~\eqref{eq:critical-norms}.

\begin{lemma}[cutoff commutator]\label{lem:commutator}
For $j=X,I$ and smooth $u=(a,\phi)$,
\begin{equation}\label{eq:commutator-identity}
 [\mathscr D_{A_\lambda,\Phi_\lambda},\beta_j]u
 =\left(-\iota_{\nabla\beta_j}a,
 (d\beta_j\wedge a)^+,\rho(d\beta_j)\phi\right),
\end{equation}
and
\begin{equation}\label{eq:commutator-estimate}
 \|[\mathscr D_{A_\lambda,\Phi_\lambda},\beta_j]u\|_{
 \mathbb Z\oplus\mathbb Y}
 \leq C(\log N)^{-1/2}\|u\|_{\mathbb X_{q_\lambda}}.
\end{equation}
The estimate extends by density to the weak domain.
\end{lemma}

\begin{proof}
All zero-order terms commute with a scalar cutoff.  The three principal
parts give~\eqref{eq:commutator-identity}.  The gauge row, the $L^2$ part of
the curvature row, and the Dirac row are bounded using
$L^4\cdot L^4\subset L^2$ and the first estimate in
\eqref{eq:log-cutoff}.  For the critical part of the curvature row,
\eqref{eq:critical-products} gives
\[
 \|(d\beta_j\wedge a)^+\|_{L^\sharp}
 \leq C\|d\beta_j\|_{L^{2\sharp}}
       \|a\|_{L^{2\sharp}}
 \leq C(\log N)^{-1/2}\|a\|_{L^2_{1,A_\lambda}}.
\]
Adding the rows proves~\eqref{eq:commutator-estimate}.  Notice that no
derivative of $d\beta_j$ occurs: the augmented operator is first order.
\end{proof}

\subsection{The induced bubble operator}
Let $I$ be the centered, framed charge-one instanton on $S^4$, and let
$K_I=\coker D_I$, a complex line.  The stabilized round operator is
\begin{equation}\label{eq:round-bubble-operator}
 \widehat B_0(a,\phi,\kappa)
 =\left(d_I^*a,d_I^+a,D_I\phi-\iota_I\kappa\right).
\end{equation}
It is inverted on the fixed $L^2_1(S^4)$ complement of translations,
dilation, and framed rotations.

A direct comparison of the physical operator with~\eqref{eq:round-bubble-operator}
would compare their principal symbols in $L^\sharp$.  Such a comparison is
not bounded merely by an $L^2_1$ domain norm.  We instead compactify the
physical normal-coordinate metric and its Clifford structure to a smooth
metric $g_{x,\lambda}$ on $S^4$, equal to the transported physical geometry
on the active bubble region and to the round geometry near the north pole.
Let $\widehat B_{x,\lambda}$ denote the operator defined by this geometry and
the full instanton $I$.

In normal coordinates $z$ at $x$, uniformly over $X$,
\begin{equation}\label{eq:normal-coordinate-metric}
 g_{ij}(z)=\delta_{ij}+O(|z|^2),\qquad
 \partial g_{ij}(z)=O(|z|),\qquad \partial^2g_{ij}(z)=O(1).
\end{equation}
Put $z=\lambda y$ and $\Omega(y)=2/(1+|y|^2)$.  On the active region
$|y|\leq\tfrac12\lambda^{-1/2}$, set
\[
 \widetilde g_{x,\lambda}=\Omega^2g_x(\lambda y).
\]
In the inverted coordinate $w=y/|y|^2$, its boundary has radius
$2\sqrt\lambda$.  Interpolate there to the round metric
$g_0=\Omega^2\delta$ with a cutoff depending on $|w|/\sqrt\lambda$.
The tensor being cut off is $O(\lambda)$, while one cutoff derivative costs
$O(\lambda^{-1/2})$.  Hence
\begin{equation}\label{eq:induced-metric-estimate}
 \|\widetilde g_{x,\lambda}-g_0\|_{C^0(g_0)}\leq C\lambda,
 \qquad
 \|\nabla^{g_0}(\widetilde g_{x,\lambda}-g_0)\|_{C^0(g_0)}
 \leq C\sqrt\lambda.
\end{equation}
The interpolation lies where $\beta_I=0$.

\begin{proposition}[induced bubble inverse]\label{prop:bubble-inverse}
After fixing the stabilizing map $\iota_I$, the operators
$\widehat B_{x,\lambda}$ have right inverses $P^I_{x,\lambda}$ on the
fixed model-kernel complement such that
\begin{align}\label{eq:induced-comparison}
 \|\widehat B_{x,\lambda}-\widehat B_0\|_{
 \mathcal L(L^2_1\oplus K_I,L^2)}&\leq C\sqrt\lambda,\\
 \|P^I_{x,\lambda}y\|_{L^2_1\oplus K_I}&\leq C\|y\|_{L^2}.
 \label{eq:induced-inverse}
\end{align}
If $I_\lambda^{\mathrm{tr}}$ is the truncated instanton used in the
pregluing, its replacement of $I$ changes the transported operator by at
most $C_N\sqrt\lambda$ from the weak domain to the mixed target.
\end{proposition}

\begin{proof}
Transport the physical zero-spinor ASD and Dirac blocks separately to
$(S^4,\widetilde g_{x,\lambda})$.  For the spinor block the conformal law is
\[
 D_A^{h^{-2}g}=h^{5/2}D_A^g h^{-3/2};
\]
the ASD block has its four-dimensional conformal transport, and the Coulomb
row is transported with the gauge-fixed ASD operator.  On
$\operatorname{supp}\beta_I$ this definition is an exact conjugation of the
physical operator.  Thus its principal symbol agrees there exactly.
Estimate~\eqref{eq:induced-metric-estimate} controls the principal
coefficients in $L^2$ and the spin-connection and gauge-row zero-order
coefficients, proving~\eqref{eq:induced-comparison}.  The Neumann formula
\[
 P^I_{x,\lambda}=P^I_0
 \bigl(1+(\widehat B_{x,\lambda}-\widehat B_0)P^I_0\bigr)^{-1}
\]
then proves~\eqref{eq:induced-inverse}.  It applies to mixed target data
because that norm contains the $L^2$ norm.  On the active region the induced
and physical principal symbols agree exactly, so no principal-symbol error
is estimated in $L^\sharp$.

In the south-pole framing, $|I(y)|\leq C|y|^{-3}$ and
$|\nabla I(y)|\leq C|y|^{-4}$.  Truncation begins at
$R\asymp\lambda^{-1/2}$.  If $b_\lambda=I_\lambda^{\mathrm{tr}}-I$, radial
integration gives
\begin{equation}\label{eq:instanton-tail}
 \|b_\lambda\|_{L^4}+\|\nabla b_\lambda\|_{L^2}
 \leq C\lambda,
 \qquad
 \|b_\lambda\|_{L^2}+\|b_\lambda\|_{L^{2\sharp}}
 \leq C\sqrt\lambda.
\end{equation}
Changing a connection changes all three rows only by multiplication by
$b_\lambda$.  The ordinary rows follow from $L^4\cdot L^4\subset L^2$;
the critical curvature row follows from~\eqref{eq:critical-products}.  This
proves the last assertion.
\end{proof}

\subsection{Localized coefficients and matched obstruction spaces}
The remaining comparison terms are zero-order.  We repeatedly use the
following two consequences of~\eqref{eq:critical-norms}.  If $f=(v,\psi)$
is a target section supported in $B(x,r)$ and is pointwise bounded by $M$,
then
\begin{equation}\label{eq:shrinking-support}
 \|v\|_{L^{\sharp,2}}+\|\psi\|_{L^2}\leq CMr^2.
\end{equation}
Indeed, the $L^2$ bounds follow from volume, while the integral of
$\operatorname{dist}(z,\cdot)^{-2}$ over a four-ball is $O(r^2)$,
uniformly in $z$.

If instead $f$ is a bounded coefficient supported in the same ball, then
\begin{equation}\label{eq:supported-coefficient}
 \|f\|_{L^4}+\|f\|_{L^{2\sharp}}
 \leq Cr\|f\|_{L^\infty},
 \qquad
 \|fu\|_{L^2}+\|fu\|_{L^\sharp}
 \leq Cr\|f\|_{L^\infty}\|u\|_{L^2_1}.
\end{equation}
The first inequality follows by volume comparison and by splitting the
supremum in~\eqref{eq:critical-norms} according as its base point lies within
$2r$ of $x$; the second follows from~\eqref{eq:critical-products} and
$L^4\cdot L^4\subset L^2$.

The pregluing formulas have the exact support properties
\begin{align}\label{eq:pregluing-supports}
 \Phi_\lambda&=\Phi_0
 &&\text{outside }B(x,8\lambda^{1/3}),
 &\Phi_\lambda&=0
 &&\text{inside }B(x,4\lambda^{1/3}),\notag\\
 A_\lambda&=A_0
 &&\text{outside }B(x,4\sqrt\lambda).
\end{align}
Set $s_\lambda=\Phi_\lambda-\Phi_0$.  The compact background family and
\eqref{eq:supported-coefficient} give
\begin{equation}\label{eq:spinor-coefficient}
 \|s_\lambda\|_{L^4}+\|s_\lambda\|_{L^{2\sharp}}
 \leq C\lambda^{1/3}.
\end{equation}
This controls the spinor coefficients in the gauge, curvature, and Dirac
rows, using $L^4$ for the ordinary rows and $L^{2\sharp}$ for the critical
curvature row.

In radial gauge, $|A_0-\Gamma|\leq C|z|$.  On
$\operatorname{supp}\beta_X$, where $|z|\geq N^{-1}\sqrt\lambda$, the
south-pole instanton decay gives
\[
 |I_\lambda(z)|\leq C\lambda^2|z|^{-3}
 \leq C N^3\sqrt\lambda.
\]
Thus $c_{x,\lambda}=\beta_X(A_\lambda-A_0)$ satisfies
\begin{equation}\label{eq:background-connection-coefficient}
 \|c_{x,\lambda}\|_{L^\infty}\leq C_N\sqrt\lambda,
 \qquad
 \|c_{x,\lambda}\|_{L^4}
 +\|c_{x,\lambda}\|_{L^{2\sharp}}\leq C_N\lambda.
\end{equation}
It contributes $O_N(\sqrt\lambda)$ to the three rows, including the Coulomb
row.  Finally, the fixed perturbing one-form in the bubble Dirac row is
supported by $\beta_I$ in a ball of radius $\tfrac12\sqrt\lambda$; ordinary
H\"older gives an $O(\sqrt\lambda)$ multiplier bound.  Consequently the
total localized coefficient error is
\begin{equation}\label{eq:coefficient-error}
 \|E_{\mathrm{spin}}+E^X_{\mathrm{conn}}+E^X_{\mathrm{gauge}}\|
 \leq C_N\lambda^{1/3}.
\end{equation}

Choose the metrics on $\Xi_q$ and $K_I$ so that the limiting embeddings
$\iota_{s,q}$ and $\iota_I$ are isometries.  With the same cutoffs as in
\eqref{eq:square-partition}, define
\begin{equation}\label{eq:matched-embeddings}
 J_{s,\lambda}\xi=\beta_XT^Y_{X,\lambda}\iota_{s,q}\xi,
 \qquad
 J_{i,\lambda}\kappa
 =T^-_{x,\lambda}\bigl(\beta_I^\#\iota_I\kappa\bigr),
\end{equation}
where $\beta_I^\#(y)=\beta_I(\lambda y)$.  The transports include the
framing, dilation, and the negative-spinor conformal weight.

\begin{lemma}[matched obstruction embeddings]\label{lem:matched-obstruction}
For all sufficiently small scales, $J_{s,\lambda}\oplus J_{i,\lambda}$ is
injective with a uniformly bounded inverse on its image.  Its $L^2$
projection is uniformly bounded in the mixed target norm.  In the patched
parametrix the two finite-dimensional stabilization terms cancel exactly.
\end{lemma}

\begin{proof}
The background section removed by $1-\beta_X$ is smooth and supported in a
ball of radius $O(\sqrt\lambda)$, so~\eqref{eq:shrinking-support} shows that
its mixed-target norm tends to zero.  On the bubble, the cutoff complement
moves to the north pole of $S^4$; compactness of the unit sphere in the
finite-dimensional space $K_I$ makes the $L^2$ tail uniformly small.  The
off-diagonal inner products tend to zero because the two supports separate
after rescaling.  The Gram matrix therefore tends uniformly to the identity.
This proves the first two assertions.  The last one follows directly from
using precisely the profiles in~\eqref{eq:matched-embeddings} both when the
model target is localized and when its obstruction coefficient is spliced
back.  Thus there is no separate obstruction error in the parametrix.
\end{proof}

\subsection{A complement which does not degenerate}
Let
\[
 \cK_X(q)=\ker\widehat{\mathscr D}^X_q,\qquad
 \cK_I=\ker\widehat B_0.
\]
The first is identified by the background finite-dimensional reduction with
$T_qN(\varepsilon)$ and has real dimension
$d_{\mathfrak s}+2r_N$: along the zero section its tangential summand is
$T_pM_{\mathfrak s}$ and its normal summand is $N_p$.  The second is the
eight-dimensional framed ASD kernel.  The physical $L^2$ norm of a naturally
dilated bubble one-form is $\lambda$ times its model norm.  We therefore do
not use physical $L^2$-orthogonality to define the complement.

The parameter metric is
\begin{equation}\label{eq:parameter-metric}
 g_{\mathrm{par},\lambda}=g_{\overline N}+\lambda^{-2}g_X
 +(d\lambda/\lambda)^2+g_{\SO(3)}.
\end{equation}
Thus the normalized directions are $\lambda\partial_{x^j}$,
$\lambda\partial_\lambda$, and $\partial_{\theta^k}$.  They converge after
rescaling to a fixed basis of $\cK_I$.

Define the synthesis map
\begin{equation}\label{eq:synthesis}
 S_{\lambda,N}(k_X,k_I)=
 \left(\beta_XT^+_{X,\lambda}u_X
 + T^+_{I,\lambda}(\beta_I^\#u_I),\eta_X,0\right),
\end{equation}
where $k_X=(u_X,\eta_X)$ and $k_I=(u_I,0)$.  Conversely, localize a physical
domain element by $\beta_X$ and $\beta_I$, transport it to the two fixed
models, project there onto the configuration components of $\cK_X(q)$ and
$\cK_I$, and lift back to the augmented kernels.  Denote the resulting
retraction by $R_{\lambda,N}$.  It ignores the two freely varying obstruction
coefficients.

Let $P^X_q$ denote the smoothly varying inverse of the stabilized background
operator on the complement of $\cK_X(q)$.  We use this inverse, the kernel
projection, and $\iota_{s,q}$ at the actual point $q$, rather than transport
them to finitely many reference points.  All their strong and mixed graph
norms are uniformly bounded on the compact disk bundle.  The background
model equations therefore cancel exactly before the neck error is
estimated.

\begin{proposition}[model-coordinate complement]\label{prop:complement}
For fixed $N$ and sufficiently small $\lambda$,
\begin{equation}\label{eq:model-matrix}
 M_{\lambda,N}=R_{\lambda,N}\operatorname{pr}_{\mathbb X}
 S_{\lambda,N}=I+o(1)
\end{equation}
uniformly in all gluing data.  Hence
\begin{equation}\label{eq:model-projection}
 \Pi_{\lambda,N}=S_{\lambda,N}M_{\lambda,N}^{-1}
 R_{\lambda,N}\operatorname{pr}_{\mathbb X}
\end{equation}
is a uniformly bounded projection onto the
$d_{\mathfrak s}+2r_N+8$ parameter modes.
Moreover,
\begin{equation}\label{eq:approximate-kernel}
 \|\widehat{\mathscr D}_\lambda S_{\lambda,N}\|
 \leq C(\log N)^{-1/2}+C_N\lambda^{1/3}.
\end{equation}
\end{proposition}

\begin{proof}
On a background kernel vector, the part removed by the cutoff is supported
in a shrinking ball.  The derivative term is bounded by
\[
 \|d\beta_Xu_X\|_{L^2}
 \leq C\|d\beta_X\|_{L^2}\|u_X\|_{L^\infty}
 \leq C_N\sqrt\lambda\|k_X\|.
\]
On a bubble kernel vector, the normalized cutoff complement escapes to the
north pole.  Compactness of the unit sphere in the fixed finite-dimensional
space gives uniform decay of its $L^2_1$ and $L^4$ tails.  The two
off-diagonal entries tend to zero: a smooth background mode acquires a
positive power of $\lambda$ under bubble rescaling, while only the escaping
tail of a bubble mode reaches the background projection.  This proves
\eqref{eq:model-matrix}; it also gives the uniform bound in
\eqref{eq:model-projection}.

Apply the global augmented operator to~\eqref{eq:synthesis}.  The two model
kernel equations remove the interior terms, Lemma~\ref{lem:commutator}
controls the cutoffs, Proposition~\ref{prop:bubble-inverse} and
\eqref{eq:coefficient-error} control the model comparisons, and
Lemma~\ref{lem:matched-obstruction} cancels the stabilization terms.  Their
sum is~\eqref{eq:approximate-kernel}.
\end{proof}

\medskip
Set
\begin{equation}\label{eq:complement-space}
 \mathbb H^{\mathrm{cmp}}_{\lambda,N}
 =\ker\bigl(R_{\lambda,N}\operatorname{pr}_{\mathbb X}\bigr).
\end{equation}
Literal normalized derivatives of the pregluing map differ by $o(1)$ in
$L^2_1$ from the synthesized modes, so they define the same uniformly
transverse complement after decreasing the scale bound.

\subsection{The exact inverse}
Let $P^I_{x,\lambda}$ be the inverse from
Proposition~\ref{prop:bubble-inverse}.  We now give the raw parametrix
explicitly.  For a smooth physical target $y$, define
\begin{align}\label{eq:target-localizations}
 L^Y_{X,\lambda}y
 &= (T^Y_{X,\lambda})^{-1}(\beta_Xy),
 &L^Y_{I,\lambda}y
 &= (T^Y_{I,\lambda})^{-1}(\beta_Iy),
\end{align}
where the second transport includes dilation and the conformal spinor weight.
The corresponding domain synthesis maps are
\begin{align}\label{eq:domain-synthesis}
 L^X_{X,\lambda}(u,\xi)
 &=\bigl(\beta_XT^+_{X,\lambda}u,\xi,0\bigr),\notag\\
 L^X_{I,\lambda}(v,\kappa)
 &=\bigl(T^+_{I,\lambda}(\beta_I^\#v),0,\kappa\bigr).
\end{align}
Thus
\begin{equation}\label{eq:raw-parametrix}
 P^{\mathrm{raw}}_{\lambda,N}
 =L^X_{X,\lambda}P^X_qL^Y_{X,\lambda}
  +L^X_{I,\lambda}P^I_{x,\lambda}L^Y_{I,\lambda}.
\end{equation}
Multiplication by the scalar cutoffs is bounded on the ordinary $L^2$ target
rows.  On the curvature row it is bounded on both $L^2$ and $L^\sharp$;
the latter follows immediately from the defining positive-kernel integral in
\eqref{eq:critical-norms}.  The transports in
\eqref{eq:target-localizations} are uniformly bounded because the background
geometry ranges over a compact family and the bubble transport is the exact
conformal transport used to define $\widehat B_{x,\lambda}$ on the support of
$\beta_I$.  Equations~\eqref{eq:target-localizations}--\eqref{eq:raw-parametrix}
therefore define a uniformly bounded map from smooth mixed-target data to the
augmented weak domain, and then by density to the graph closure.  The
parametrix used below is $(1-\Pi_{\lambda,N})P^{\mathrm{raw}}_{\lambda,N}$.

An arbitrary $L^2_1$ correction need not have curvature row in
$L^{\sharp,2}$.  We therefore define $\mathbb X^{\mathrm{gr}}_{q_\lambda}$
to be the completion of smooth corrections and auxiliary coefficients in
the norm
\begin{align}\label{eq:graph-domain}
 \|(u,\eta_s,\eta_i)\|_{\mathrm{gr}}
 ={}&\|u\|_{L^2_1}+\|\eta_s\|+\|\eta_i\|\notag\\
 &+\|\widehat{\mathscr D}_\lambda(u,\eta_s,\eta_i)\|_{
       \mathbb Z\oplus\mathbb Y}.
\end{align}
The augmented operator is bounded from this graph domain by definition.
Elliptic regularity identifies every graph-domain solution of the homogeneous
equation with a strong Sobolev solution.

\begin{theorem}[uniform augmented inverse]\label{thm:uniform-inverse}
There is $C>0$ such that for every sufficiently large $N$ there are
$\varepsilon,\delta_N>0$ for which
\begin{equation}\label{eq:parametrix-bound}
 \left\|\widehat{\mathscr D}_\lambda
 (1-\Pi_{\lambda,N})P^{\mathrm{raw}}_{\lambda,N}-I\right\|
 \leq C(\log N)^{-1/2}+C_N\lambda^{1/3}.
\end{equation}
For $q\in\overline N(\varepsilon)$ and $0<\lambda\leq\delta_N$, the restriction
\begin{equation}\label{eq:restricted-operator}
 \widehat{\mathscr D}_\lambda:
 \mathbb H^{\mathrm{cmp}}_{\lambda,N}\cap
 \mathbb X^{\mathrm{gr}}_{q_\lambda}\longrightarrow
 \mathbb Z_{q_\lambda}\oplus\mathbb Y_{q_\lambda}
\end{equation}
is an isomorphism on the strong Fredholm level.  Its inverse
$P_{\lambda,N}$ obeys
\begin{equation}\label{eq:uniform-inverse-bound}
 \|P_{\lambda,N}y\|_{\mathrm{gr}}
 \leq C\|y\|_{\mathbb Z\oplus\mathbb Y}.
\end{equation}
The inverse is smooth and equivariant on every positive-scale stratum.
\end{theorem}

\begin{proof}
Apply $\widehat{\mathscr D}_\lambda$ to the two summands in
\eqref{eq:raw-parametrix}.  Their principal terms are $\beta_X^2y$ and
$\beta_I^2y$, which recombine to $y$ by
\eqref{eq:square-partition}.  Lemma~\ref{lem:commutator} contributes
$C(\log N)^{-1/2}$; Proposition~\ref{prop:bubble-inverse},
\eqref{eq:coefficient-error}, and variation of the background operator
at the actual point $q$ contribute
$C_N\lambda^{1/3}$; the matched
obstruction error is zero.  Projecting the output adds at most the same
quantity by~\eqref{eq:approximate-kernel}.  This proves
\eqref{eq:parametrix-bound}.

Choose $N$ first so that the logarithmic term is less than $1/4$, and then
choose $\delta_N$ so that the scale term is less than $1/4$.  A Neumann
series gives a right inverse with image in
\eqref{eq:complement-space} and proves~\eqref{eq:uniform-inverse-bound}.
The strong augmented operator has index $d_{\mathfrak s}+2r_N+8$.  The
complement has codimension $d_{\mathfrak s}+2r_N+8$, because
\eqref{eq:model-matrix} is invertible.  The
restricted operator has index zero; surjectivity therefore implies
injectivity.  Elliptic regularity identifies the weakly constructed inverse
with this strong inverse.  The estimate extends to the mixed-target graph
closure.  Every ingredient is smooth and equivariant at positive scale, and
inversion in the open group of bounded invertible operators preserves these
properties.
\end{proof}

\begin{lemma}[normalized derivatives of the inverse]
\label{lem:inverse-derivatives}
Trivialize the mixed targets by the background and splicing transports and
the complements by $1-\Pi_{\lambda,N}$.  For $j=0,1,2$, covariant
differentiation with respect to the normalized parameter metric
\eqref{eq:parameter-metric} gives
\begin{equation}\label{eq:inverse-derivative-bound}
 \|\nabla^jP_{\lambda,N}\|_{
 \mathcal L(\mathbb Z\oplus\mathbb Y,
            \mathbb X^{\mathrm{gr}}_{q_\lambda})}\leq C_j.
\end{equation}
The same statement holds for the strong inverse in strong Sobolev norms.
\end{lemma}

\begin{proof}
Write $Q_\lambda=(1-\Pi_{\lambda,N})P^{\mathrm{raw}}_{\lambda,N}$ and
$E_\lambda=\widehat{\mathscr D}_\lambda Q_\lambda-I$.  One or two
normalized derivatives preserve the parametrix estimate.  Background
derivatives range over a compact unit sphere bundle;
$\lambda\partial_{x^j}$ and $\lambda\partial_\lambda$ conjugate under
bubble dilation to fixed derivatives of the unit-scale model; frame
derivatives act in fixed compact-group representations; and differentiated
cutoff commutators obey the same logarithmic bound.  The same estimates hold
for synthesis, retraction, and $\Pi_{\lambda,N}$.  Consequently
\[
 \|\nabla^jE_\lambda\|
 \leq C_j(\log N)^{-1/2}+C_{j,N}\lambda^{1/3},
 \qquad 0\leq j\leq2.
\]
Since
$P_{\lambda,N}=Q_\lambda(I+E_\lambda)^{-1}$, differentiate this identity
and use
\[
 \nabla(I+E)^{-1}=-(I+E)^{-1}(\nabla E)(I+E)^{-1}.
\]
The second derivative contains the corresponding $\nabla^2E$ term and two
ordered products of $\nabla E$.  This proves the graph-domain and strong
bounds.
\end{proof}

In particular, a normalized sequence in the complement cannot be an
approximate kernel supported in the neck: applying
\eqref{eq:uniform-inverse-bound} would give $1\leq o(1)$.  This also excludes
gauge modes.  On the background the nonzero spinor removes the
determinant-one stabilizer, and on the bubble it is removed by the based
framing convention.

\section{Projected residual and nonlinear correction}\label{sec:correction}

\subsection{The seven-region estimate}
Put $r=\sqrt\lambda$ and $R=\lambda^{1/3}$.  For small $\lambda$ the
pregluing formula divides $X$ into the core $B(x,r/4)$; the inner, flat, and
middle annuli with radii $(r/4,r/2)$, $(r/2,2r)$, and $(2r,4r)$; the
zero-spinor plateau $(4r,4R)$; the outer spinor-cutoff annulus $(4R,8R)$; and
the exterior of $B(x,8R)$.  The logarithmic parametrix cutoffs lie inside the
$r$-scale annuli but play no role in this residual decomposition.

\begin{proposition}[projected residual]\label{prop:residual}
Uniformly in the compact background family, the center, and the frames,
\begin{equation}\label{eq:residual}
 \|\mathfrak S(q_\lambda)-J_{s,\lambda}b(q)\|_{L^{\sharp,2;2}}
 \leq C\lambda^{1/3}.
\end{equation}
Here $L^{\sharp,2;2}$ means $L^{\sharp,2}$ in the curvature component and
$L^2$ in the negative-spinor component.
\end{proposition}

\begin{proof}
On the exterior, the preglued pair equals the background pair and
$J_{s,\lambda}$ equals the background obstruction embedding; the residual
therefore vanishes identically.

Write $\mathfrak S(q)=(v_0,\psi_0)$ and let $\chi_{\mathrm{out}}$ be the outer
spinor cutoff.  On $4R<d(x,\cdot)<8R$ the pair is
$(A_0,\chi_{\mathrm{out}}\Phi_0)$ and
$J_{s,\lambda}b(q)=\mathfrak S(q)$.  If $\mathcal Q(\Phi)$ denotes the
quadratic spinor term in the curvature equation, the residual is
\begin{equation}\label{eq:outer-residual}
 \bigl(\mathcal Q(\Phi_0)-\mathcal Q(\chi_{\mathrm{out}}\Phi_0),
 (\chi_{\mathrm{out}}-1)\psi_0
 +\rho(d\chi_{\mathrm{out}})\Phi_0\bigr).
\end{equation}
The first component and the first term in the second component are uniformly
bounded on a ball of radius $8R$, and hence have mixed-target norm $O(R^2)$
by~\eqref{eq:shrinking-support}.  Direct radial integration gives
$\|d\chi_{\mathrm{out}}\|_{L^2}\leq CR$; consequently the last term in
\eqref{eq:outer-residual} is $O(R)=O(\lambda^{1/3})$.

On the plateau $4r<d(x,\cdot)<4R$ the pair is $(A_0,0)$, not a solution of
the background Dirac equation.  Here the exact difference from the matched
background residual is
\begin{equation}\label{eq:plateau-residual}
 \mathfrak S(A_0,0)-\mathfrak S(A_0,\Phi_0)
 =\bigl(\mathcal Q(\Phi_0),-\psi_0\bigr),
\end{equation}
with the sign in $\mathcal Q$ fixed by the convention in
\eqref{eq:outer-residual}.  Both entries are uniformly bounded and supported
in $B(x,4R)$, so their mixed-target norm is $O(R^2)=O(\lambda^{2/3})$.

On the middle annulus $2r<d(x,\cdot)<4r$, write the connection in radial
gauge as $\Gamma+\chi_{\mathrm{mid}}(A_0-\Gamma)$.  Then
\begin{equation}\label{eq:middle-curvature}
 F=\chi_{\mathrm{mid}}F_{A_0}
   +d\chi_{\mathrm{mid}}\wedge(A_0-\Gamma)
   +(\chi_{\mathrm{mid}}^2-\chi_{\mathrm{mid}})
    (A_0-\Gamma)\wedge(A_0-\Gamma).
\end{equation}
The bounds $|A_0-\Gamma|\leq Cr$ and
$|d\chi_{\mathrm{mid}}|\leq C/r$ show that
\eqref{eq:middle-curvature} is uniformly bounded.  The spinor is zero and the
matched background target is uniformly bounded, so
\eqref{eq:shrinking-support} gives an $O(r^2)=O(\lambda)$ contribution.  On
the flat annulus $r/2<d(x,\cdot)<2r$ the preglued pair is $(\Gamma,0)$; the
only term after subtraction is again a uniformly bounded matched background
target, and its contribution is $O(\lambda)$.

On the inner annulus $r/4<d(x,\cdot)<r/2$, the instanton decay gives
\begin{equation}\label{eq:inner-decay}
 |F_I|\leq C\lambda^2d(x,\cdot)^{-4}\leq C,
 \qquad
 |A_I|\leq C\lambda^2d(x,\cdot)^{-3}\leq C\sqrt\lambda.
\end{equation}
Since the cutoff derivative is $O(\lambda^{-1/2})$, every curvature term in
the truncated connection and every subtracted background term is uniformly
bounded.  This annulus therefore also contributes $O(\lambda)$.

Finally, on $B(x,r/4)$ the model instanton is anti-self-dual for the Euclidean
metric $\delta$, but not identically for the physical metric $g$.  Its
self-dual curvature is
\begin{equation}\label{eq:core-metric-error}
 F_I^{+,g}=\tfrac12(*_g-*_\delta)F_I.
\end{equation}
The normal-coordinate estimate $|g-\delta|\leq Cd(x,\cdot)$, conformal
invariance of the four-dimensional curvature norm, and radial integration of
the standard instanton give
\begin{equation}\label{eq:core-metric-bound}
 \|F_I^{+,g}\|_{L^{\sharp,2}(B(x,r/4))}\leq C\sqrt\lambda.
\end{equation}
The matched background target on the core contributes $O(r^2)=O(\lambda)$.
Adding the seven disjoint regional estimates gives
\[
 C\bigl(\lambda^{1/3}+\lambda^{2/3}
       +\lambda+\lambda^{1/2}\bigr)
 \leq C\lambda^{1/3},
\]
which proves~\eqref{eq:residual}.
\end{proof}

\subsection{Solving the augmented equation}
Write $h=(u,\eta_s,\eta_i)$ and let $P_\lambda$ be the inverse in
Theorem~\ref{thm:uniform-inverse}.  Expanding~\eqref{eq:augmented-map} gives
\begin{equation}\label{eq:nonlinear-expansion}
 \widehat\cF_\lambda(h)=r_\lambda+
 \widehat{\mathscr D}_\lambda h+Q_\lambda(u,u),
\end{equation}
where $r_\lambda$ is the residual in Proposition~\ref{prop:residual}.
Sobolev multiplication gives the scale-independent estimates
\begin{align}\label{eq:quadratic-estimates}
 \|Q_\lambda(u,u)\|_{\mathbb Z\oplus\mathbb Y}
 &\leq C\|u\|_{L^2_1}^2,\notag\\
 \|Q_\lambda(u,u)-Q_\lambda(v,v)\|_{\mathbb Z\oplus\mathbb Y}
 &\leq C(\|u\|_{L^2_1}+\|v\|_{L^2_1})\|u-v\|_{L^2_1}.
\end{align}

\begin{proposition}[small correction]\label{prop:correction}
For sufficiently small $\lambda$, there is a unique
$h_\lambda=(u_\lambda,\eta_{s,\lambda},\eta_{i,\lambda})$ in a fixed small
ball in $\mathbb H^{\mathrm{cmp}}_{\lambda,N}$ satisfying
$\widehat\cF_\lambda(h_\lambda)=0$.  It obeys
\begin{equation}\label{eq:correction}
 \|u_\lambda\|_{L^2_1}+\|\eta_{s,\lambda}\|
 +\|\eta_{i,\lambda}\|\leq C\lambda^{1/3}.
\end{equation}
It is smooth in all positive-scale gluing parameters.
\end{proposition}

\begin{proof}
On the ball of radius $2C\|r_\lambda\|$, use the map
\[
 h\longmapsto-P_\lambda\bigl(r_\lambda+Q_\lambda(u,u)\bigr).
\]
Theorem~\ref{thm:uniform-inverse}, Proposition~\ref{prop:residual}, and
\eqref{eq:quadratic-estimates} show that it preserves this ball and has
Lipschitz constant less than $1/2$ after decreasing the scale.  Banach's
fixed-point theorem gives existence, uniqueness, and~\eqref{eq:correction}.
The strong Sobolev inverse and the implicit-function theorem give smooth
dependence for $\lambda>0$.
\end{proof}

Define the corrected gluing map and obstruction section by
\begin{align}\label{eq:corrected-map-section}
 \gamma(g)&=[q_\lambda+u_\lambda],\notag\\
 \chi(g)&=(b(q)+\eta_{s,\lambda},\eta_{i,\lambda}).
\end{align}
Then $\gamma(g)$ is a genuine monopole exactly when $\chi(g)=0$.

\subsection{Continuity at zero scale}
The constants above are independent of a moving center.  If
$g_n\to g_0$ with $\lambda_n\to0$, then~\eqref{eq:correction} gives
$u_{\lambda_n}\to0$ in the global weak norm.  On compact subsets away from
the limiting center, the preglued pairs converge in every Sobolev norm to
the background.  After rescaling at the centers, their connections converge
to the framed instanton while their spinors converge to zero.  These two
statements, together with the energy identity, are precisely Uhlenbeck
convergence to $(q,x)$.  Thus $\gamma$ extends continuously by
\begin{equation}\label{eq:zero-scale-map}
 \gamma(q,x,0)=(q,x).
\end{equation}
The same estimates and the definitions in~\eqref{eq:matched-embeddings}
show that the norm of $\chi$ is continuous when the instanton line is
collapsed at the cone vertex.

\section{Parameter recovery and the compactified embedding}
\label{sec:embedding}

\subsection{Uniform modulation}
Work first on the framed cover of the positive-scale gluing domain.  At a
preglued pair $q_\lambda(g)$ impose the based Coulomb condition
\begin{equation}\label{eq:based-slice}
 d^{0,*}_{q_\lambda(g)}\bigl((A,\Phi)-q_\lambda(g)\bigr)=0.
\end{equation}
For $q_g=(A_g,\Phi_g)$, put
\begin{equation}\label{eq:pair-laplacian}
 \Delta_g^0=d_g^{0,*}d_g^0
 =d_{A_g}^*d_{A_g}+\mathfrak r_{\Phi_g}^*\mathfrak r_{\Phi_g}.
\end{equation}
Let $L^{\sharp,2}_{2,g}$ be the completion of smooth determinant-one
infinitesimal gauges in the norm
\begin{equation}\label{eq:critical-gauge-norm}
 \|\zeta\|_{2,\sharp,g}
 =\|\Delta_g^0\zeta\|_{L^{\sharp,2}}+\|\zeta\|_{L^2}.
\end{equation}
Feehan's critical embedding gives
$L^{\sharp,2}_{2,g}\hookrightarrow C^0\cap L^2_2$
\cite[Lemmas~5.3--5.6]{FeehanCritical}.

\begin{theorem}[uniform critical slice for pairs]
\label{thm:critical-pair-slice}
After decreasing the maximum scale, there is $C>0$ such that
\begin{equation}\label{eq:critical-green-estimate}
 \|\zeta\|_{2,\sharp,g}
 \leq C\|\Delta_g^0\zeta\|_{L^{\sharp,2}}
\end{equation}
for every positive-scale family datum.  The Green operator
$G_g=(\Delta_g^0)^{-1}$ and its first two normalized parameter derivatives
are uniformly bounded.  Moreover, there is $r_{\mathrm{sl}}>0$, independent
of the datum, such that every nearby based gauge orbit has a unique
representative satisfying~\eqref{eq:based-slice}.
\end{theorem}

\begin{proof}
We first prove a uniform ordinary spectral gap.  Otherwise there are data
$g_j$ and infinitesimal gauges $\zeta_j$ with $\|\zeta_j\|_{L^2}=1$ and
\[
 \|d_{A_{g_j}}\zeta_j\|_{L^2}
 +\|\zeta_j\Phi_{g_j}\|_{L^2}\longrightarrow0.
\]
Kato's inequality and Sobolev embedding give a common $L^4$ bound.  Away
from the bubbling point, Rellich compactness gives a limit that is parallel
and annihilates the nonzero limiting spinor; hence it vanishes.  No $L^2$
mass can remain at the point, since
\[
 \|\zeta_j\|_{L^2(B(x_j,2r))}
 \leq Cr\|\zeta_j\|_{L^4(B(x_j,2r))}\leq Cr.
\]
This contradicts the normalization and proves the gap.

For $f=\Delta_g^0\zeta$, positivity of
$\mathfrak r_{\Phi_g}^*\mathfrak r_{\Phi_g}$ and the weak Kato inequality
give
\[
 (d^*d+1)|\zeta|\leq |f|+|\zeta|
\]
in distributions.  The scalar Green-kernel estimate and the spectral gap
therefore yield
\[
 \|\zeta\|_{C^0}+\|\zeta\|_{L^2}
 \leq C\|f\|_{L^{\sharp,2}}.
\]
Applying the covariant-Laplacian estimates of
\cite[Lemmas~5.5--5.6]{FeehanCritical} to
$d_{A_g}^*d_{A_g}\zeta
=f-\mathfrak r_{\Phi_g}^*\mathfrak r_{\Phi_g}\zeta$
proves~\eqref{eq:critical-green-estimate}.  Direct differentiation of
\eqref{eq:pair-laplacian}, followed by the critical multiplication
estimates and
$\nabla G=-G(\nabla\Delta^0)G$, proves the first two normalized derivative
bounds.  The normalization in~\eqref{eq:parameter-metric} cancels the
inverse lengths of the center and scale profiles.

The contraction argument for the gauge equation gives existence of a
Coulomb representative with a uniform radius.  For injectivity, suppose
$q_g+c_1$ and $q_g+c_2$ are in the slice and related by a based gauge
$u$; set $v=u-1$.  The transition equations and the two Coulomb equations
give the exact identity
\begin{equation}\label{eq:nonlinear-slice-identity}
 \Delta_g^0v
 =d_{A_g}^*(va_1-a_2v)-\mathfrak r_{\Phi_g}^*(v\phi_1).
\end{equation}
Pairing with $v$, using the ordinary gap just proved and H\"older's
inequality, gives
\[
 \|d_{A_g}v\|_{L^2}^2+\|v\Phi_g\|_{L^2}^2
 \leq C(\|c_1\|_{L^2_1}+\|c_2\|_{L^2_1})
 \bigl(\|d_{A_g}v\|_{L^2}^2+\|v\Phi_g\|_{L^2}^2\bigr).
\]
Decrease $r_{\mathrm{sl}}$ and absorb.  Then $v=0$ by the gap, and the two
representatives agree.  This also excludes a slowly rotating transition
along the long neck.  Smooth dependence follows from the same contraction;
weak continuity follows from the uniform bound and uniqueness.
\end{proof}

The Coulomb-horizontal projection is
\begin{equation}\label{eq:horizontal-projection}
 H_g=1-d_g^0G_gd_g^{0,*}.
\end{equation}
It and its first two normalized derivatives are uniformly bounded by
Theorem~\ref{thm:critical-pair-slice}.

Put
\begin{equation}\label{eq:horizontal-complement}
 R_g=R_{\lambda,N}\operatorname{pr}_{\mathbb X},
 \qquad
 \cC_g^{\mathrm{sl}}=\ker R_g\cap\ker d^{0,*}_{q_\lambda(g)}.
\end{equation}
After locally trivializing this bundle, define the quotient-valued tubular map
\begin{equation}\label{eq:tubular-map}
 \mathcal T(g,h)=[q_\lambda(g)+h],
 \qquad h\in\cC_g^{\mathrm{sl}}.
\end{equation}
Its differential along the zero section is
\begin{equation}\label{eq:modulation-differential}
 D\mathcal T_{(g,0)}(\dot g,\dot h)
 =H_gdq_\lambda(\dot g)+\dot h.
\end{equation}
Let
$\jmath_g:T_g\cG^+\to\cK_X(q)\oplus\cK_I$ be the coordinate identification
which sends the normalized parameter basis from
\eqref{eq:parameter-metric} to the background virtual-normal modes and the
translation, dilation, and framed-rotation Jacobi fields.  Differentiating the
splicing formulas gives
\begin{equation}\label{eq:parameter-retraction}
 \|R_gH_gdq_\lambda-\jmath_g\|
 \longrightarrow0
 \quad\text{uniformly as }\lambda\downarrow0.
\end{equation}
For background variations, the discrepancy in
\eqref{eq:parameter-retraction} is supported on a shrinking ball.  For
$\lambda\partial_x$, $\lambda\partial_\lambda$, and
$\partial_\theta$, rescaling gives the fixed instanton Jacobi fields.
Derivatives of an $r=\sqrt\lambda$ cutoff contribute the compensating factors
$\lambda/r=\sqrt\lambda$ and $\lambda^2/r^2=\lambda$ for one and two
normalized center derivatives; logarithmic-scale derivatives preserve the
fixed cutoff profiles.  The remaining terms are smooth background terms on
shrinking annuli or escaping instanton tails.  The estimates used in
Proposition~\ref{prop:complement} apply to each such derivative.  Horizontal
projection changes the result by $o(1)$ because the gauge row has the uniform
coercive bound just established.  This proves
\eqref{eq:parameter-retraction}.

Since $R_g\dot h=0$, applying $R_g$ to
\eqref{eq:modulation-differential} first recovers $\dot g$ with a uniformly
bounded inverse; substitution then recovers $\dot h$.  Thus
$D\mathcal T_{(g,0)}$ is uniformly invertible.  The preceding differentiated
cutoff estimates also give a uniform bound for the normalized second
derivative of $\mathcal T$.  The quantitative Banach inverse-function theorem,
in the fixed local trivializations, consequently gives constants
$c_{\mathrm{mod}},r_{\mathrm{mod}}>0$ independent of sufficiently small scale.

\begin{lemma}[uniform modulation]\label{lem:modulation}
The map~\eqref{eq:tubular-map} is one-to-one on
\begin{equation}\label{eq:modulation-tube}
 d_{\mathrm{par}}(g,g')<c_{\mathrm{mod}},
 \qquad
 \|h\|_{L^2_1}+\|h'\|_{L^2_1}<r_{\mathrm{mod}}.
\end{equation}
Equivalently, after the unique based gauge in~\eqref{eq:based-slice} is
imposed, every class in this tube has a unique parameter $g$ characterized by
\begin{equation}\label{eq:modulation-characterization}
 R_g\bigl((A,\Phi)-q_\lambda(g)\bigr)=0.
\end{equation}
The recovered parameter and based gauge are continuous in the weak topology
and smooth on the positive-scale strong Sobolev charts.  The construction is
equivariant for the diagonal frame actions and therefore descends to the
gluing quotient.
\end{lemma}

The correction in Proposition~\ref{prop:correction} satisfies both equations
in~\eqref{eq:horizontal-complement}.  Hence exact modulation recovers the
input parameter of every corrected configuration: no post-correction shift of
center, scale, or relative frame is hidden in the construction.

\subsection{Center, scale, and frame as observables}
For an exact framed finite-energy connection on $\RR^4$, the model center
and scale are
\begin{equation}\label{eq:model-center-scale}
 z[A]=\frac{1}{8\pi^2}\int_{\RR^4}y\,|F_A|^2,
 \qquad
 \sigma[A]^2=\frac{1}{8\pi^2}
 \int_{\RR^4}|y-z[A]|^2|F_A|^2.
\end{equation}
They identify translations and dilations of the charge-one instanton.  For
the corrected configurations we use only bounded local observables: a fixed
fractional-energy radius, a first moment truncated in normalized
coordinates, and convergence of the rescaled framed connection on compact
sets.  Proposition~\ref{prop:correction} implies that these observables tend
to those of the input instanton.  We do not apply the unbounded global
moments in~\eqref{eq:model-center-scale} directly to the corrected
connection.

\begin{lemma}[outer-normalized bubble framing]
\label{lem:outer-frame}
Let $q_{g_n}+u_n$ be corrected configurations with
$\lambda_n\downarrow0$ and
\[
 \|u_n\|_{L^2_1}+\|\eta_{s,n}\|+\|\eta_{i,n}\|
 \leq C\lambda_n^{1/3}.
\]
Normalize the exterior gauges by the fixed based background slice.  The
rescaled gauges may then be chosen to converge on compact subsets of
$S^4-\{\infty\}$.  In temporal gauge the constant cylinder mode determines
a value at bubble infinity, modulo the center, and this value is precisely
the relative frame.  The resulting outer-normalized framed bubble limit is
unique.
\end{lemma}

\begin{proof}
Use the dyadic Coulomb gauges and the stabilized long-neck estimate of
Proposition~\ref{prop:stabilized-long-neck} below.  Starting with the identity
at the exterior end, remove the nonconstant temporal-gauge modes while
retaining the constant endpoint $\rho_n\in\SO(3)$.  If $k_j$ relates the
remaining gauges on successive slabs, the fixed-annulus estimate gives
\[
 \operatorname{dist}(k_{j+1},k_j)
 \leq C(\alpha_j+\alpha_{j+1}+\beta_j+\beta_{j+1}),
\]
where $\alpha_j$ and $\beta_j$ are the connection and spinor norms on that
slab.  The two-sided cylinder estimate gives a length-independent bound for
the sum of these quantities, tending to zero after the exterior and interior
ends are sent out.  Thus every nonconstant change of gauge on the neck
vanishes in the limit, while compactness of $\SO(3)$ gives a limit of the
retained endpoints.  Fixed-annulus elliptic estimates and the transition
equation $dv_n=v_na_n-a_n'v_n$ give convergence of the relating gauges on
compact sets.  Two outer-normalized limits differ by a stabilizer of the
irreducible instanton with the same value at infinity, hence only by the
central action already divided out.
\end{proof}

\begin{proposition}[separation of parameters]\label{prop:separation}
Suppose $g_n$ and $g'_n$ are gluing parameters tending to the compactified
domain and $\gamma(g_n)=\gamma(g'_n)$.  Then their background points and
ideal centers agree in the limit.  If the scales are positive, their scale
ratios tend to one and their relative frames tend to the same element of
$\SO(3)$.  For all sufficiently small scale, $g_n=g'_n$.
\end{proposition}

\begin{proof}
If the scales stay bounded away from zero, the smooth local inverse in
Lemma~\ref{lem:modulation} gives the result.  It remains to consider
$\lambda_n,\lambda'_n\to0$.  Uhlenbeck convergence first identifies the two
background limits and ideal points.  The fixed background slices then give
$d_N(q_n,q'_n)\to0$.

Rescale the common corrected configuration about $x_n$ by $\lambda_n$.
The first representation converges on compact sets to the centered
unit-scale instanton $I$.  In these coordinates, the second bubble has
relative translation and scale
\begin{equation}\label{eq:relative-center-scale}
 z_n=\lambda_n^{-1}\exp_{x_n}^{-1}(x'_n),
 \qquad t_n=\lambda'_n/\lambda_n.
\end{equation}
If $|z_n|\to\infty$, the second description has no bubble energy on a fixed
compact set.  If $t_n\to\infty$, it becomes flat there; if $t_n\to0$, its
curvature measure tends to a point mass.  Each alternative contradicts the
smooth unit-scale limit.  Thus, after a subsequence,
$z_n\to z\in\RR^4$ and $t_n\to t\in(0,\infty)$.

The two rescaled limits are gauge equivalent, and the second is the
translate and dilate of $I$ by $(z,t)$.  Applying
\eqref{eq:model-center-scale} to these exact limiting instantons gives
$z=0$ and $t=1$.  Apply Lemma~\ref{lem:outer-frame} to both descriptions,
using the same exterior based normalization.  Their retained values at
bubble infinity coincide, so their relative $\SO(3)$ frames agree modulo
exactly the central and diagonal actions already divided out.  Consequently
\begin{equation}\label{eq:parameter-separation}
 d_{\mathrm{par}}(g_n,g'_n)\longrightarrow0.
\end{equation}
For large $n$, the two parameters lie within the uniform modulation radius.
Put the common configuration in the based slice and apply
Lemma~\ref{lem:modulation}; uniqueness gives literal equality.  At zero
scale the cone quotient has already identified all relative frames.
\end{proof}

Because the compactified parameter space over the compact background is
compact after the outer boundary is fixed, sequential separation also rules
out coincidences between parameters that are not initially close.  We have
therefore proved the following statement.

\begin{corollary}[compactified embedding]\label{cor:embedding}
After decreasing $\varepsilon$ and $\delta$, $\gamma$ is an equivariant
topological embedding of the compactified virtual gluing space.  On the
positive-scale stratum it is a smooth embedding.
\end{corollary}

Finally write $\gamma(g)=q_g+u_g$ in the chosen based slice.  For
$0\leq t\leq1$, set
\begin{equation}\label{eq:isotopy}
 \gamma_t(g)=q_g+t\,u_g.
\end{equation}
The parameter-separation argument is uniform in $t$, because
$\|t u_g\|_{L^2_1}\leq C\lambda^{1/3}$ and the derivative of the modulation
map remains $M_{\lambda,N}+o(1)$.  Hence every $\gamma_t$ is a compactified
embedding.  The construction is equivariant, $\gamma_0$ is the splicing
map, and $\gamma_1=\gamma$.

\section{Reverse gluing and neighborhood surjectivity}
\label{sec:reverse-gluing}

Continuity and injectivity do not imply that every nearby monopole is
obtained by gluing.  We prove the converse by first putting an arbitrary
one-bubble sequence close to a preglued pair and then applying the same
complement and inverse used in the construction.

\begin{proposition}[strong splicing approximation]
\label{prop:strong-splicing}
Suppose that irreducible monopoles $[A_n,\Phi_n]$ converge in the Uhlenbeck
topology to
\[
 ([B,\Psi],x)\in M_{\mathfrak s}\times X
\]
with exactly one unit of instanton charge lost at $x$.  There are coarse
gluing data
\[
 g_n^0=([B,\Psi],x_n,\lambda_n,\rho_n),\qquad
 x_n\to x,\quad\lambda_n\downarrow0,
\]
standard gluing bundle identifications, and determinant-one gauges for which
\begin{equation}\label{eq:strong-splicing}
 \|(A_n,\Phi_n)-q_{\lambda_n}(g_n^0)\|_{L^2_1}\longrightarrow0.
\end{equation}
The relative frames $\rho_n$ use the quotient convention in
\eqref{eq:gluing-quotient}; no auxiliary tangent or bundle frame remains in
the datum.
\end{proposition}

\subsection{Extraction of the bubble}
The monopole Weitzenb\"ock formula gives a uniform $L^\infty$ bound for
$\Phi_n$ on the fixed compact perturbation family.  Choose $x_n$ and
$\lambda_n$ by a fixed fractional-energy centering and radius convention in
a normal-coordinate ball about $x$.  Pull back by
$y\mapsto\exp_{x_n}(\lambda_n y)$ and choose Uhlenbeck gauges.  The rescaled
curvature equation has the form
\begin{equation}\label{eq:rescaled-curvature}
 F_{A_n^{\mathrm{res}}}^+
 =\lambda_n^2\mathcal Q(\Phi_n^{\mathrm{res}})
  +\lambda_n^2\mathcal P_n,
\end{equation}
where $\mathcal P_n$ is bounded on each fixed normalized ball.  The right
side tends to zero.  Uhlenbeck compactness and removal of the singularity at
infinity therefore give a finite-energy ASD connection $I$ on $S^4$.  The
lost charge is one, so $I$ has charge one.  The centering and radius
conventions make it centered and of unit scale; the framed charge-one
classification identifies the remaining parameter with a based rotation
$\rho_n\in\SO(3)$.  We choose this rotation by comparing the rescaled bubble
gauge with the outer background gauge through the annular patching
normalization below; this is equivalent to the framed classification and
prevents a second, unrecorded frame from entering the gluing datum.

After a subsequence,
\begin{align}\label{eq:two-sided-strong-convergence}
 A_n^{\mathrm{res}}&\longrightarrow I
 &&\text{strongly in $L^2_1$ on compact subsets of $\RR^4$,}\notag\\
 (A_n,\Phi_n)&\longrightarrow(B,\Psi)
 &&\text{strongly in $L^2_1$ on compact subsets of $X\setminus\{x\}$.}
\end{align}
The conformally rescaled spinor tends to zero on the bubble side.

On a conformal cylinder, the connection gauge-curvature block has tangential
operator
\begin{equation}\label{eq:odd-signature-operator}
 L_0(\xi,b)=(-d_{S^3}^*b,-d_{S^3}\xi+*_{S^3}d_{S^3}b),
 \qquad L=L_0\oplus D_{S^3}.
\end{equation}
For the unit round sphere, the exact-form spectrum of $L_0$ is
$\{\pm\sqrt{k(k+2)}:k\geq1\}$, its coexact spectrum is
$\{\pm(k+1):k\geq1\}$, and its kernel consists exactly of the constant
$\xi$ modes.  The spin Dirac spectrum is
$\{\pm(k+3/2):k\geq0\}$.  Hence, after removal of the constant gauge mode,
\begin{equation}\label{eq:cylinder-gap}
 \operatorname{spec}L\cap(-3/2,3/2)=\varnothing.
\end{equation}

\begin{lemma}[fixed-annulus coupled estimate]
\label{lem:fixed-annulus}
Let $D'\Subset D$ be one of the fixed round annulus pairs obtained by
rescaling a dyadic annulus.  Let $(A,\Phi)$ solve the monopole equation on
$D$, and let $(C,\Psi)$ solve it up to a defect $f\in L^p$, for a fixed
$2<p<4$.  Suppose both connections are in Coulomb gauge relative to the
flat product connection and have scale-invariant
$L^4\oplus W^{1,2}$ norm below the Uhlenbeck constant.  Then, for
$z=(A-C,\Phi-\Psi)$,
\begin{equation}\label{eq:fixed-annulus-estimate}
 \|z\|_{W^{1,p}(D')}
 \leq C_p\bigl(\|z\|_{L^2(D)}+\|f\|_{L^p(D)}\bigr).
\end{equation}
On an overlap, the based mismatch $w$ between the actual and approximate
transition maps satisfies the corresponding
$W^{1,4}\cap W^{2,2}\cap W^{2,p}$ estimate with right-hand side given by
the two neighboring $z$-norms and the defect.
\end{lemma}

\begin{proof}
Subtract the gauge, curvature, and Dirac equations.  The result is a fixed
first-order elliptic operator applied to $z$, with coefficients given by the
two connection matrices and spinors, equal to $f$.  Multiplication by a
small $L^4$ connection coefficient maps $W^{1,p}$ to $L^p$, because
\[
 W^{1,p}\hookrightarrow L^{4p/(4-p)},\qquad
 \frac1p=\frac14+\frac{4-p}{4p},
\]
and can be absorbed.  Interior elliptic regularity proves
\eqref{eq:fixed-annulus-estimate}.  Subtracting the two transition equations
first gives the $W^{1,4}$ estimate for $w-1$; differentiating once and using
\eqref{eq:fixed-annulus-estimate} gives the $W^{2,2}$ and $W^{2,p}$ bounds on
a smaller annulus.  Since $p>2$, the latter embeds in $C^0$, placing $w$ in
one fixed logarithm chart.
\end{proof}

\begin{lemma}[two-sided cylinder estimate]
\label{lem:two-sided-cylinder}
Suppose on $[j_-,j_+]\times S^3$ that $P_0z=0$ and
\begin{equation}\label{eq:cylinder-equation}
 (\partial_t+L)z=R(t)z+Q(z)+f,
\end{equation}
where the operator norm of $R$ and the Lipschitz norm of $Q$ on each
two-slab enlargement are sufficiently small.  Let $\alpha_j$ be the
scale-invariant $W^{1,2}\cap W^{1,p}$ norm of $z$ there, for a fixed
$2<p<4$, and let $\beta_j$ be the corresponding $L^2\cap L^p$ norm of
$f$.  There are $C>0$ and $0<\theta<1$, independent of the cylinder length,
such that
\begin{align}\label{eq:three-annulus}
 \alpha_j\leq C\biggl(&\theta^{j-j_-}\alpha_{j_-}
 +\theta^{j_+-j}\alpha_{j_+}
 +\sum_{k=j_-}^{j_+}\theta^{|j-k|}\beta_k\biggr),\\
 \sum_{j=j_-}^{j_+}\alpha_j
 &\leq C\left(\alpha_{j_-}+\alpha_{j_+}
 +\sum_{j=j_-}^{j_+}\beta_j\right).
 \label{eq:cylinder-l1}
\end{align}
\end{lemma}

\begin{proof}
Let $P_+$ and $P_-$ be the positive and negative spectral projections of
$L$.  Variation of constants, propagated from the left for $P_+$ and from
the right for $P_-$, gives
\begin{align*}
 P_+z(t)&=e^{-(t-a)L}P_+z(a)
 +\int_a^te^{-(t-s)L}P_+h(s)\,ds,\\
 P_-z(t)&=e^{-(t-b)L}P_-z(b)
 -\int_t^be^{-(t-s)L}P_-h(s)\,ds,
\end{align*}
where $h=Rz+Q(z)+f$.  The gap~\eqref{eq:cylinder-gap} bounds the semigroups
by $e^{-3|t-s|/2}$.  Interior elliptic estimates give
\eqref{eq:three-annulus}.  The terms involving $Rz+Q(z)$ are the same
geometric convolution with a small multiple of $\alpha_k$; its operator is
bounded on both $\ell^1$ and $\ell^\infty$ and can be absorbed.  Summing the
remaining geometric series proves~\eqref{eq:cylinder-l1}.
\end{proof}

\begin{proposition}[family long-neck estimate]
\label{prop:stabilized-long-neck}
Let $(A_n,\Phi_n)$ be either genuine monopoles compared with their coarse
pregluings, or corrected configurations satisfying the full stabilized
equation
\begin{equation}\label{eq:stabilized-neck-equation}
 \mathfrak S(q_{g_n}+u_n)
 =J_{s,g_n}(b(q_n)+\eta_{s,n})+J_{i,g_n}\eta_{i,n},
 \qquad d_{q_{g_n}}^{0,*}u_n=0,
\end{equation}
with
$\|u_n\|_{L^2_1}+\|\eta_{s,n}\|+\|\eta_{i,n}\|
\leq C\lambda_n^{1/3}$.  Normalize at the exterior end.  On
$\Omega_n(d,R)=B(x_n,d)-B(x_n,R\lambda_n)$ there is one determinant-one
gauge in which
\begin{align}\label{eq:family-long-neck}
 &\|A_n-A_n^{\mathrm{spl}}\|_{L^4}
 +\|\nabla_{A_n^{\mathrm{spl}}}(A_n-A_n^{\mathrm{spl}})\|_{L^2}
 +\|d_{A_n^{\mathrm{spl}}}^*(A_n-A_n^{\mathrm{spl}})\|_{L^{\sharp,2}}
 \notag\\
 &\qquad\leq C\bigl(\alpha_{n,-}+\alpha_{n,+}
 +d^{3/2}+d^2+R^{-2}+\lambda_n^{1/3}\bigr).
\end{align}
For genuine monopoles the last term may be omitted.  The constant is
independent of the number of dyadic annuli and of the point in the compact
background family.  The constant cylinder mode is retained at the inner end
as the relative frame.
\end{proposition}

\begin{proof}
Put $r=e^{-t}$ and conformally change the Euclidean neck metric to
$dt^2+g_{S^3}$.  A Euclidean spinor becomes
$\Phi_{\mathrm{cyl}}=r^{3/2}\Phi_{\mathrm E}$.  After radial and Coulomb
gauge, the relative equation has principal part $\partial_t+L$.  The
normal-coordinate metric error, bounded background curvature, quadratic
spinor term, and fixed perturbation have cylindrical size $O(r^2)$; the
charge-one tail has size $O((\lambda_n/r)^2)$.  On the at most two slabs
meeting the outer spinor cutoff, conformal covariance gives the additional
source
\[
 c(dt)(\partial_t\chi_{\mathrm{out}})r^{3/2}\Phi_{\mathrm E},
\]
of size $O(r^{3/2})$.  Thus the pregluing defect on a slab of radius $r_j$
satisfies
\begin{equation}\label{eq:neck-defect}
 \beta_j\leq C\left(r_j^2+\frac{\lambda_n^2}{r_j^2}
 +\mathbf1_{\{|j-j_{\mathrm{out}}|\leq1\}}r_j^{3/2}\right),
\end{equation}
and its sum is bounded by $C(d^{3/2}+d^2+R^{-2})$.

For~\eqref{eq:stabilized-neck-equation}, the $b(q_n)$ terms cancel on
subtraction.  The background obstruction fields contribute
$O(r_j^2+r_j^{5/2})$, while the normalized charge-one adjoint-kernel spinor
contributes $O(\lambda_n/r_j)$.  Multiplication by the coefficient bound
$C\lambda_n^{1/3}$ and geometric summation add at most
$C\lambda_n^{1/3}$ to~\eqref{eq:cylinder-l1}.

No-neck energy makes the coefficients in~\eqref{eq:cylinder-equation}
uniformly small.  Solve temporal gauge from the exterior end, removing the
constant mode from $z$ but retaining its inner endpoint.  Lemma
\ref{lem:two-sided-cylinder} gives the required length-independent
$\ell^1$ estimate.  On neighboring fixed annuli, Lemma
\ref{lem:fixed-annulus} gives, for the based transition mismatch $w_j$,
\[
 \|w_j-1\|_{W^{1,4}\cap W^{2,2}\cap W^{2,p}}
 \leq C(\alpha_j+\alpha_{j+1}+\beta_j+\beta_{j+1}).
\]
Because $p>2$, these maps lie in one logarithm chart.  Extend their
logarithms with fixed cutoffs, first over even and then over odd overlaps.
The $\ell^1$ estimate controls the critical divergence norm without a neck
length factor and proves~\eqref{eq:family-long-neck}.  The constant cocycle
propagates uniquely from the exterior; its inner endpoint is the declared
relative frame rather than a mode to be divided out.
\end{proof}

\begin{lemma}[compatible gauges on a long neck]
\label{lem:neck-gauge-patching}
Let
\[
 \Omega_n(d,R)=B(x_n,d)\setminus B(x_n,R\lambda_n),
\]
and let $A_n$ and $C_n$ be determinant-one connections on the same splicing
bundle over this annulus.  Suppose that their total curvature energies on
$\Omega_n(d,R)$ tend to zero in the iterated order $n\to\infty$,
$d\downarrow0$, and $R\to\infty$.  Suppose also that, on fixed-ratio collars
at the two ends, chosen gauges for $A_n$ and $C_n$ converge respectively to
the same background connection and, after the declared rotation $\rho_n$, to
the same framed bubble connection.  Then there is one determinant-one gauge
$u_n$ on the whole neck, agreeing with these endpoint normalizations on
smaller collars, such that
\begin{align}\label{eq:patched-neck-estimate}
 &\|u_n(A_n)-C_n\|_{L^4(\Omega_n(d,R))}
 +\|\nabla_{C_n}(u_n(A_n)-C_n)\|_{L^2(\Omega_n(d,R))}\notag\\
 &\quad\leq C\bigl(
 \|F_{A_n}\|_{L^2(\Omega_n(d,R))}
 +\|F_{C_n}\|_{L^2(\Omega_n(d,R))}\bigr)
 +o_{n,d,R}(1),
\end{align}
where $C$ is independent of the number of dyadic annuli and the last term is
the sum of the two endpoint discrepancies.
\end{lemma}

\begin{proof}
Use enlarged dyadic annuli $D_j$ and concentric cores $D'_j$ so that the
$D'_j$ cover the neck and the enlarged family has overlap bounded by an
absolute constant.  After rescaling each $D_j$ to a fixed annulus, Uhlenbeck's
small-curvature theorem~\cite{Uhlenbeck82} gives based Coulomb gauges
$a_j$ and $c_j$ for $A_n$ and $C_n$ with
\begin{equation}\label{eq:local-neck-gauges}
 \|a_j\|_{L^4(D_j)}+\|\nabla a_j\|_{L^2(D_j)}
 +\|c_j\|_{L^4(D_j)}+\|\nabla c_j\|_{L^2(D_j)}
 \leq C e_j,
\end{equation}
where
\[
 e_j^2=\|F_{A_n}\|_{L^2(D_j)}^2
       +\|F_{C_n}\|_{L^2(D_j)}^2.
\]
Bounded overlap gives
\begin{equation}\label{eq:neck-square-sum}
 \sum_j e_j^2\leq C\int_{\Omega_n^+}
        (|F_{A_n}|^2+|F_{C_n}|^2),
\end{equation}
where $\Omega_n^+$ enlarges only the two endpoint collars.

We record the overlap step because it is what prevents a logarithmic loss.
If $v_j$ is the transition map between two neighboring Coulomb gauges, then
on the fixed-shape overlap
\begin{equation}\label{eq:transition-equation}
 dv_j=v_ja_{j+1}-a_jv_j.
\end{equation}
Poincare's inequality applied to~\eqref{eq:transition-equation}, followed by
the polar projection of the average to the compact determinant-one group,
produces a constant $k_j$ such that
\begin{equation}\label{eq:transition-estimate}
 \|k_j^{-1}v_j-1\|_{W^{1,4}}
 +\|k_j^{-1}v_j-1\|_{W^{2,2}}
 \leq C(e_j+e_{j+1}).
\end{equation}
The $W^{2,2}$ estimate follows by differentiating
\eqref{eq:transition-equation}; its right-hand side is in $L^2$ by
$L^4\cdot L^4\subset L^2$, and the Coulomb equations control the weak first
derivatives of $a_j$ and $a_{j+1}$.  The same construction gives constants
$\ell_j$ for the transitions of the $c_j$ gauges.  On a slightly smaller
overlap, the gauge-fixed monopole equations for $(a_j,\Phi_n)$ and the
explicit equation for the smooth spliced pair give, by interior elliptic
regularity, a bound in $W^{2,p}$ for the transition maps for some fixed
$p>2$.  Its norm relative to the constants is bounded by a modulus
$\omega(e_j+e_{j+1})$ with $\omega(t)\to0$.  Hence
$W^{2,p}\hookrightarrow C^0$ puts each based transition in one normal
neighborhood of the identity.  The group logarithm and cutoff extension used
in the next paragraph are therefore well-defined.  This bootstrap is used
only to construct the gauges; the global quantitative estimate remains the
critical square-sum estimate~\eqref{eq:neck-square-sum}.

Fix the outer background normalization.  Since the nerve of the annular cover
is an interval, multiply successive local gauges by constants recursively so
that the constant parts $k_j$ and $\ell_j$ agree.  The remaining overlap maps
are the based maps in~\eqref{eq:transition-estimate}.  The standard cutoff
extension on the fixed overlap absorbs each such map into one of its two
neighbors.  Its $W^{1,4}\cap W^{2,2}$ norm is bounded by the right-hand side of
\eqref{eq:transition-estimate}.  Applying the extensions on alternating
overlaps first, and then on the remaining overlaps, gives global gauges.  The
resulting $L^4\oplus W^{1,2}$ errors are estimated by the square sum, not by
their ordinary sum.  Thus~\eqref{eq:neck-square-sum} gives a constant
independent of the length of the neck.

It remains to identify the normalization at the inner end.  Cap the outer end
by the background bundle and the inner end by the charge-one bubble bundle.
The two capped bundles have the same second Chern class: their only clutching
class is the unit charge already used in the splicing construction.  Hence the
product of the constant overlap transitions contains no additional
$\pi_3(\SO(3))$ class.  After the outer constant has been fixed, its inner
constant is precisely the constant comparing the rescaled inner gauge with
the centered framed instanton; by our convention following
\eqref{eq:rescaled-curvature}, this constant is $\rho_n$.  Inserting
$\rho_n$ into $C_n$ makes the endpoint constants agree.  The endpoint
convergence absorbs the two final collar errors into $o_{n,d,R}(1)$ and proves
\eqref{eq:patched-neck-estimate}.
\end{proof}

\begin{lemma}[one-bubble no neck]\label{lem:no-neck}
For the centers and scales above,
\begin{equation}\label{eq:no-neck-energy}
 \lim_{R\to\infty}\lim_{d\downarrow0}\limsup_{n\to\infty}
 \int_{B(x_n,d)\setminus B(x_n,R\lambda_n)}|F_{A_n}|^2=0.
\end{equation}
After the bundle identifications determined by $\rho_n$, there are gauges
on these necks in which the actual connection and the spliced
background--instanton connection satisfy
\begin{equation}\label{eq:no-neck-strong}
 \|A_n-A_n^{\mathrm{spl}}\|_{L^4}
 +\|\nabla_{A_n^{\mathrm{spl}}}(A_n-A_n^{\mathrm{spl}})\|_{L^2}
 \longrightarrow0
\end{equation}
in the iterated order $n\to\infty$, $d\downarrow0$, and $R\to\infty$.
\end{lemma}

\begin{proof}
The curvature measures converge to
\[
 |F_B|^2dV+8\pi^2\delta_x.
\]
For fixed $d$, the region outside $B(x_n,d)$ accounts for the background
energy.  After rescaling, $B(x_n,R\lambda_n)$ accounts for the energy of $I$
on $B(0,R)$.  Taking the three limits in the displayed order, the background
energy in the shrinking ball and the instanton tail both vanish.  There is
no second concentration point, so~\eqref{eq:no-neck-energy} follows.

Cover the neck by dyadic annuli.  Once their curvature energy is below the
Uhlenbeck threshold, the general small-curvature Coulomb-gauge estimate
\cite{Uhlenbeck82} gives
\begin{equation}\label{eq:annular-gauge}
 \|a\|_{L^4}+\|\nabla a\|_{L^2}\leq C\|F_A\|_{L^2}.
\end{equation}
This is the estimate for a general connection, not an ASD annular decay
theorem.  It applies here because the full neck energy tends to zero.  The
monopole equation and the uniform spinor bound also give
\begin{equation}\label{eq:self-dual-small-ball}
 \|F_{A_n}^+\|_{L^2(B(x_n,d))}\leq Cd^2.
\end{equation}
The annuli have bounded overlap, so the squares of
\eqref{eq:annular-gauge} sum to a constant times the total neck energy; no
factor equal to the number of annuli occurs.  Apply the same estimate to the
spliced connection, whose neck energy is small by bounded background
curvature and instanton-tail decay.
The two-sided convergence in~\eqref{eq:two-sided-strong-convergence} supplies
the endpoint hypotheses of Lemma~\ref{lem:neck-gauge-patching}; the spliced
connection uses the same outer background gauge and the same inner rotation
$\rho_n$.  Apply~\eqref{eq:patched-neck-estimate} and then the three limits.
This proves~\eqref{eq:no-neck-strong} in one global determinant-one gauge.
\end{proof}

\begin{proof}[Proof of Proposition~\ref{prop:strong-splicing}]
Given $\epsilon>0$, first choose $R$ so that the $L^2_1$ tail of $I$ outside
$B(0,R)$ is smaller than $\epsilon$.  Next choose $d$ so that the background
has small local energy in $B(x,2d)$ and the no-neck quantity is smaller than
$\epsilon$.  For all large $n$, the first convergence in
\eqref{eq:two-sided-strong-convergence} controls $B(x_n,R\lambda_n)$,
Lemma~\ref{lem:no-neck} controls the intermediate annulus, and the second
convergence controls $X\setminus B(x_n,d)$.

The spinor has no missing neck term.  Its pointwise bound gives
\begin{equation}\label{eq:spinor-small-ball}
 \|\Phi_n\|_{L^4(B(x_n,2d))}=O(d),
\end{equation}
and the Dirac equation with an interior covariant Caccioppoli estimate gives
\begin{equation}\label{eq:spinor-caccioppoli}
 \|\nabla_{A_n}\Phi_n\|_{L^2(B(x_n,d))}\leq Cd.
\end{equation}
Indeed, the Weitzenb\"ock formula bounds the square of the left side by the
cutoff contribution, bounded zeroth-order terms, and
\[
 \int_{B(x_n,2d)}|F_{A_n}|\,|\Phi_n|^2
 \leq \|\Phi_n\|_{L^\infty}^2\|F_{A_n}\|_{L^2}
       \operatorname{vol}(B(x_n,2d))^{1/2}=O(d^2).
\]
The cutoff background spinor satisfies the same estimates.  Patching the
three comparisons with the pregluing profiles yields
\[
 \limsup_{n\to\infty}
 \|(A_n,\Phi_n)-q_{\lambda_n}(g_n^0)\|_{L^2_1}\leq C\epsilon.
\]
Letting $\epsilon$ tend to zero proves~\eqref{eq:strong-splicing}.
\end{proof}

\subsection{Exact modulation and complemented uniqueness}
Apply Lemma~\ref{lem:modulation} to the configurations in
Proposition~\ref{prop:strong-splicing}.  After based determinant-one gauge,
there are unique exact gluing data
$g_n=(q_n,\widehat x_n,\widehat\lambda_n,\widehat\rho_n)$ and differences
$u_n$ such that
\begin{equation}\label{eq:exact-modulation}
 (A_n,\Phi_n)=q_{\widehat\lambda_n}(g_n)+u_n,\qquad
 d^{0,*}_{q_{\widehat\lambda_n}}u_n=0,\qquad
 R_{g_n}u_n=0,\qquad
 \|u_n\|_{L^2_1}\to0.
\end{equation}
The exact parameters are close to the coarse ones in
\eqref{eq:parameter-metric}; in particular, $q_n$ tends to the reducible
background and $b(q_n)\to0$.

Set
\begin{equation}\label{eq:reverse-augmented-vector}
 H_n=(u_n,-b(q_n),0).
\end{equation}
The complement condition depends only on $u_n$, so $H_n$ lies in
$\mathbb H^{\mathrm{cmp}}_{\widehat\lambda_n,N}$.  Since $(A_n,\Phi_n)$ is
a monopole and is in the based slice,
\begin{equation}\label{eq:reverse-augmented-equation}
 \widehat\cF_{g_n}(H_n)=0.
\end{equation}
Let $e_n=\widehat\cF_{g_n}(0)$.  The inverse and quadratic estimates give
\[
 \|H_n\|
 \leq C\|e_n\|+C\|H_n\|^2
 \leq C\widehat\lambda_n^{1/3}+C\|H_n\|^2.
\]
The qualitative convergence in~\eqref{eq:exact-modulation} absorbs the
quadratic term, and hence
\begin{equation}\label{eq:reverse-rate}
 \|H_n\|\leq 2C\widehat\lambda_n^{1/3}.
\end{equation}
Thus $H_n$ lies in the same contraction ball as the correction in
Proposition~\ref{prop:correction}.  Uniqueness of that fixed point gives
\begin{equation}\label{eq:reverse-uniqueness}
 (u_n,-b(q_n),0)
 =(u_{g_n},\eta_{s,g_n},\eta_{i,g_n}).
\end{equation}
It follows that $[A_n,\Phi_n]=\gamma(g_n)$ and $\chi(g_n)=0$.

\begin{proposition}[neighborhood surjectivity]\label{prop:neighborhood}
There is an Uhlenbeck-open neighborhood $\cU$ of
$M_{\mathfrak s}\times X$ such that
\begin{equation}\label{eq:neighborhood-image}
 \cU\cap\overline\cM_{\mathfrak t'}^{*,0}
 =\cU\cap\gamma(\chi^{-1}(0)).
\end{equation}
Consequently, $\gamma$ restricts to a homeomorphism from $\chi^{-1}(0)$
onto an Uhlenbeck-open neighborhood of the ideal stratum, after restricting
the virtual neighborhood over its outer boundary.
\end{proposition}

\begin{proof}
The gluing construction gives the inclusion from right to left.  If the
reverse inclusion failed in every neighborhood, there would be genuine
monopoles outside $\gamma(\chi^{-1}(0))$ converging to the compact ideal
stratum.  A subsequence converges to one point of that stratum, and the
argument above places every sufficiently late member in
$\gamma(\chi^{-1}(0))$, a contradiction.  The inverse is continuous because
$\gamma$ is a compactified embedding by Corollary~\ref{cor:embedding}.
\end{proof}

\section{The obstruction cone and orientations}\label{sec:orientations}

\subsection{The compactified section}
The background summand $\overline\Upsilon^s$ is the pullback of the fixed
complex bundle $\Xi\to M_{\mathfrak s}$ to
$N_{\mathfrak t,\mathfrak s}$.  The charge-one Dirac
operator has the kernel and cokernel in~\eqref{eq:bubble-cokernel}.
At positive scale its obstruction line is
\[
 (0,\delta]\times\SU(2)\times_{\{\pm1\}}\CC
 \longrightarrow(0,\delta]\times\SO(3).
\]
Its compactification is the cone map~\eqref{eq:obstruction-cone}; the whole
fiber, not merely its norm, is identified with the vertex at zero scale.
There are no partition-overlap maps at level one.

At positive scale, define $\chi$ by~\eqref{eq:corrected-map-section}; at zero
scale, set
\begin{equation}\label{eq:zero-scale-obstruction}
 \chi(q,x,0)=(b(q),0).
\end{equation}
The contraction is equivariant, because the nonlinear map, cutoffs,
transports, and inverse are equivariant and its fixed point is unique.  We use
the diagonal presentation before taking the effective residual circle; it has
weight two on the normal and background-obstruction factors.  After passage
to the effective action, the corresponding Euler class is
\begin{equation}\label{eq:background-euler}
 e(\overline\Upsilon^s/S^1)=(-\nu)^{r_\Xi}.
\end{equation}
The instanton line retains the nontrivial
$\SU(2)/\{\pm1\}$ framing twist, so its Euler class is computed from its
tensor square in Section~\ref{sec:pairing} rather than copied from
\eqref{eq:background-euler}.

\begin{proposition}[compactified obstruction section]
\label{prop:obstruction-section}
The two components of $\chi$ are smooth at positive scale.  The background
component extends continuously over the compactification and smoothly on
each zero-scale stratum.  The instanton component is a continuous section of
the cone pseudo-bundle.  With the metric induced by
$J_{s,g}\oplus J_{i,g}$, the norm of $\chi$ is continuous along the section.
\end{proposition}

\begin{proof}
Positive-scale smoothness follows from Proposition~\ref{prop:correction}.
At zero scale, its estimate gives
\[
 \|\eta_{s,g}\|+\|\eta_{i,g}\|\leq C\lambda^{1/3}.
\]
Thus $\chi_s$ tends to $b(q)$ and $\chi_i$ tends to the cone vertex.  For
$g_n\to(q,x,0)$, the Gram estimate in
Lemma~\ref{lem:matched-obstruction} gives
\[
 \|J_{s,g_n}\chi_s(g_n)+J_{i,g_n}\chi_i(g_n)\|_{L^2}^2
 =\|b(q)\|^2+o(1),
\]
which is the asserted norm continuity.  This does not extend the ambient
instanton line as a constant-rank bundle; only the section tends to its
collapsed fiber.
\end{proof}

We next compare the global reduction with the local small-eigenvalue
reduction of~\cite[Section~8]{FeehanLenessIII}.  There are three distinct
finite-dimensional reductions.  On a relatively compact positive-scale
chart $U^+$ on which a
cutoff $\mu$ avoids the spectrum, let
\begin{equation}\label{eq:three-obstruction-models}
 E_{\mathrm{glob},g}=\Xi_q\oplus K_I,\qquad
 E_{\mathrm{loc},g}=\ker d_q^{1,*}\oplus K_I,\qquad
 E_{\mathrm{spec},g}=\operatorname{Ran}\Pi_{q_g,\mu}.
\end{equation}
The published cut-and-project and synthesis estimates
\cite[Propositions~8.19 and~8.26]{FeehanLenessIII} identify
$E_{\mathrm{loc}}$ with $E_{\mathrm{spec}}$ on a constant-rank chart:
their two composites are $I+o(1)$.  They also transport the ordered
orientation local system of
$\ker d_q^{1,*}\oplus K_I$ to $E_{\mathrm{spec}}$, compatibly on overlaps.
There is generally no isomorphism between $E_{\mathrm{glob}}$ and
$E_{\mathrm{loc}}$, since the global augmentation may contain extra
stabilizing directions.

\begin{proposition}[stable comparison of reductions]
\label{prop:stable-reduction}
On every positive-scale constant-rank germ of the monopole zero set, the
global and spectral reductions admit a common equivariant finite-dimensional
enlargement.  Their zero loci agree, the comparisons satisfy the cocycle
condition up to controlled equivariant homotopy, and the induced determinant
orientation and relative Euler--Thom pairing agree.  No equality of the two
ranks and no spectral bundle at zero scale is asserted.
\end{proposition}

\begin{proof}
For a smooth Fredholm section $F$ with two surjective augmentations
$i_j:E_j\to Y$, form
\[
 \cV_{01}=\{(x,e_0,e_1):F(x)-i_0e_0-i_1e_1=0\}.
\]
The two reductions embed by setting the other coefficient equal to zero.
At a point of the first reduction, projection to $e_1$ identifies its normal
bundle in $\cV_{01}$ with $E_1$: surjectivity of
$DF\oplus(-i_0)$ supplies a lift of every prescribed $e_1$.  The analogous
statement holds with the indices reversed.  Thus both obstruction zero loci
are canonically $F^{-1}(0)$.

Apply this construction with $E_0=E_{\mathrm{glob}}$ and
$E_1=E_{\mathrm{spec}}$.  The first augmented derivative is surjective by
Theorem~\ref{thm:uniform-inverse}.  On $E_{\mathrm{spec}}^\perp$, the second
has right inverse
\[
 d_{q_g}^{1,*}(d_{q_g}^1d_{q_g}^{1,*})^{-1},
\]
bounded by the spectral gap at $\mu$.  Lemma
\ref{lem:inverse-derivatives} and the Riesz formula for
$\Pi_{q_g,\mu}$ give controlled tubular maps on a common germ.  On overlaps,
order all spectral factors by chart label and use literal zero-coordinate
embeddings; associativity of direct sum gives the cocycle before tubular
choices, and nearby invariant tubular choices are homotopic through their
graph maps.

Along the first embedding,
\[
 \det T\cV_{01}\cong\det T\cV_0\otimes\det E_1.
\]
The same $E_1$ occurs in the common obstruction bundle, so its determinant
line cancels its dual.  Along the other embedding the fixed order contributes
the Koszul factor
$(-1)^{\rank_{\RR}E_0\rank_{\RR}E_1}$; it is $+1$ because
$E_0=\Xi\oplus K_I$ has even real rank.  If $E_1$ is not globally oriented,
its orientation local system is carried by both factors and cancels
canonically.  The identical cancellation holds in relative equivariant
cohomology: the Thom class of the added normal $E_1$ is the Euler factor of
the added obstruction $E_1$.  Hence stabilization changes neither the
oriented zero locus nor the relative Euler--Thom pairing.
\end{proof}

\subsection{The exact transversality statement}
At a positive-scale zero, the extended equation is
\begin{equation}\label{eq:extended-identity}
 \mathfrak S(\gamma(g))=J_g\chi(g).
\end{equation}
Let $Q_g$ be projection onto the obstruction image.  The contraction
parametrizes the local virtual manifold $(1-Q_g)\mathfrak S=0$, and the
complemented inverse identifies its tangent space with
$D\gamma_g(T_g\cG^+)$.  Differentiating~\eqref{eq:extended-identity} at a
zero gives
\begin{equation}\label{eq:differentiated-obstruction}
 D\mathfrak S_{\gamma(g)}D\gamma_g=J_gD\chi_g.
\end{equation}

\begin{proposition}[stratumwise transversality]\label{prop:transversality}
For a generic perturbation at which the irreducible monopole section is
transverse, the total section $\chi_s\oplus\chi_i$ is transverse on
$\cG^+$.  On
$(N_{\mathfrak t,\mathfrak s}-M_{\mathfrak s})\times X$, the surviving
background component $\chi_s=b$ is transverse.  No transversality holds or
is asserted on $M_{\mathfrak s}\times X$ when $r_\Xi>0$.
\end{proposition}

\begin{proof}
Given $e\in\Upsilon_g$ at a positive-scale zero, surjectivity of
$D\mathfrak S_{\gamma(g)}$ supplies $v$ with
$D\mathfrak S_{\gamma(g)}v=J_ge$.  Its projected equation component
vanishes, so $v$ is tangent to the virtual manifold and equals
$D\gamma_g\dot g$ for some $\dot g$.  Equation
\eqref{eq:differentiated-obstruction} and injectivity of $J_g$ give
$D\chi_g\dot g=e$.

On the middle stratum, the instanton fiber has collapsed and the background
finite-dimensional reduction makes $b$ transverse; the $X$ factor is
passive.  On the lowest stratum $b$ is identically zero.  An identically zero
section of a positive-rank bundle is not transverse.  This is why
\cite[Theorem~3.8 and Section~5.1]{FeehanLenessLevelOne} asserts precisely
the first two conclusions and disclaims the third.
\end{proof}

Accordingly, the all-strata formulation in
\cite[Hypothesis~7.8.1(2)]{FeehanLenessCobordism} must be read using the
coarser stratification there.  For the fine decomposition
\eqref{eq:three-strata}, the relative Euler construction uses Proposition
\ref{prop:transversality} and does not require transversality on the lowest
stratum.

\subsection{Determinant-line excision}
For a real Fredholm operator $F$, use
\[
 \det F=\Lambda^{\max}\ker F\otimes
 (\Lambda^{\max}\coker F)^*.
\]
Direct-sum factors are ordered as written, and complex vector spaces carry
their complex orientations.  For $q\in N(\varepsilon)$, write
$p=\pi_N(q)\in M_{\mathfrak s}$.  Orient $T_pM_{\mathfrak s}$ using the homology
orientation $\Omega$.  Orient the centered framed instanton parameters
by
\begin{equation}\label{eq:quaternionic-orientation}
 (\partial_\lambda,i,j,k)
 \quad\text{on}\quad \RR\oplus\mathfrak{so}(3)=\mathbb H.
\end{equation}

We use the real determinant-line excision principle of
\cite[Section~3]{DonaldsonOrientation}; see also
\cite[Section~7.1]{DK90}.  In the form needed here, suppose two ordered pairs
of real first-order elliptic operators are obtained from one another by
exchanging compact pieces, and that their bundles, principal symbols, and
operators agree on product collars after a homotopy of zero-order terms.
There is then a natural isomorphism
\begin{equation}\label{eq:abstract-excision}
 \operatorname{Exc}:
 \det D_1\otimes\det D_2
 \longrightarrow
 \det D'_1\otimes\det D'_2.
\end{equation}
It is continuous for families, compatible with composition and direct sum,
and is the identity when the exchanged data are identical.  These
normalizations determine its effect on orientations.  This is the excision
principle used to define
$O^{\mathrm{asd}}(\Omega,w)$ in
\cite[Sections~2.1--2.2]{FeehanLenessII}.

We apply~\eqref{eq:abstract-excision} to the augmented operators, for which
the parametrix supplies all required uniform Fredholm control.  Replace the
splicing annulus by a cylinder $[-T,T]\times S^3$ and denote the resulting
operator by $\widehat{\mathscr D}_{g,T}$.  On the central third of the
cylinder, homotope the coefficients to the product coefficients and remove
the zero-order terms.  On the two remaining thirds use the same transports
as in~\eqref{eq:raw-parametrix}.  The commutator and coefficient estimates of
Section~\ref{sec:parametrix} show that, after $T$ is large, this entire
homotopy has a uniformly invertible restriction to the model-coordinate
complement.  Thus its determinant line is transported without a spectral
crossing on that complement.

Cut at the middle of the product cylinder and cap the two pieces by the
fixed background and bubble models.  Excision gives
\begin{equation}\label{eq:augmented-excision}
 \det\widehat{\mathscr D}^{X}_q
 \otimes\det\widehat B_I
 \cong
 \det\widehat{\mathscr D}_{q_g}
 \otimes\det D_{\mathrm{ref}}.
\end{equation}
Here $D_{\mathrm{ref}}$ is the fixed based product operator obtained by
gluing the two discarded flat caps.  Choose its fixed finite-dimensional
stabilization so that it is invertible.  Its determinant is then canonically
identified with $\RR$ by $1\in\Lambda^0(0)$, and this factor is suppressed
below.  Equivalently, the isomorphism after suppressing it is the determinant
of the kernel--cokernel exact sequence constructed by the patched parametrix.
Because the three augmented operators are surjective, that exact sequence
reduces to the kernel gluing isomorphism
\begin{equation}\label{eq:kernel-excision}
 \ker\widehat{\mathscr D}_{q_g}
 \cong
 \ker\widehat{\mathscr D}^{X}_q
 \oplus\ker\widehat B_I.
\end{equation}
The synthesis and retraction maps in
\eqref{eq:synthesis}--\eqref{eq:model-matrix} represent this isomorphism; its
matrix is $I+o(1)$ and therefore has positive determinant.

The ordered model kernels and finite-dimensional stabilizations give
\begin{align}\label{eq:model-determinants}
 \det\widehat{\mathscr D}^{X}_q
 &=\det N_p\otimes\det T_pM_{\mathfrak s},\notag\\
 \det\widehat B_I
 &=\det T_IM_1^s(S^4,\delta),\\
 \det\widehat{\mathscr D}_{q_g}
 &=\det\mathscr D_{q_g}\otimes\det\Xi_p\otimes\det K_I.
 \notag
\end{align}
The last identity is the determinant exact sequence for adding the
finite-dimensional domain $\Xi_p\oplus K_I$ with the ordered map
$-(J_{s,\lambda}\oplus J_{i,\lambda})$.  Its minus sign has positive
determinant because the real stabilization dimension is $2r_\Xi+2$.
Combining~\eqref{eq:augmented-excision}--\eqref{eq:model-determinants} and
putting the obstruction determinants on the dual side yields
\begin{align}\label{eq:determinant-splitting}
 \det\mathscr D_{q_g}\cong{}&
 \det N_p\otimes\det T_pM_{\mathfrak s}
 \otimes\det T_I M_1^{s,\natural}(S^4,\delta)
 \otimes\det T_xX\notag\\
 &\otimes(\det\Xi_p)^*\otimes(\det K_I)^*.
\end{align}
Here the centered bubble factor contains scale and relative rotation; its
four translations are identified orientation-preservingly with $T_xX$ by
the chosen oriented tangent frame.  Moving the four-dimensional translation
factor past the four-dimensional scale--rotation factor contributes no sign.

The splitting~\eqref{eq:determinant-splitting} has no orientation monodromy
around a loop in $M_{\mathfrak s}$.  The tangent bundle is oriented by
$\Omega$, the bundles $N$ and $\Xi$ have their complex orientations, the
instanton modes have the fixed quaternionic orientation, $TX$ is oriented,
and $K_I$ is a complex line.  Possible spectral flow of an associated
self-adjoint family changes neither this real orientation nor the global
stabilization.

The tangent- and adjoint-frame determinants introduced before taking the
diagonal $\SO(4)\times\SO(3)$ quotient cancel between the ordered tangent
exact sequences for the framed parameter space and for the orbit directions.
Both group dimensions are moved past identical factors, so the two Koszul
signs cancel.  The remaining relative-frame factor is the $\SO(3)$ already
contained in $M_1^{s,\natural}$.

There is also no orientation monodromy around this relative-frame factor.
The generator of $\pi_1(\SO(3))$ lifts to a path in $\SU(2)$ from $1$ to
$-1$.  Equivariance of~\eqref{eq:abstract-excision} transports the determinant
along this lift.  At the endpoint, $-1$ acts trivially on adjoint connection
directions, by complex multiplication by $-1$ on each complex spinor or
obstruction factor, and by $-1$ on the four-dimensional quaternionic
parameter space.  Each real determinant is $+1$.  Hence the excision
orientation descends around the nontrivial loop in $\SO(3)$.

Finally, the path $q_g+t u_g$ transports
$\det\mathscr D_{q_g}$ to $\det\mathscr D_{\gamma(g)}$.  Intermediate
configurations need not solve the equations.  Replacing the published
obstruction profiles by the matched profiles is an isotopy through
embeddings, and the model-coordinate Gram map remains positive.  The
normalization discrepancy in the spinor-bilinear term is removed by a
zero-order homotopy through positive coefficients.  Thus
\eqref{eq:determinant-splitting}, with the same ordered factors, holds for the
corrected operator and is the convention of
\cite[Section~3.9]{FeehanLenessLevelOne}.

\begin{proposition}[orientation of the corrected chart]
\label{prop:chart-orientation}
At a positive-scale zero of $\chi$, there is a canonical orientation
isomorphism
\begin{equation}\label{eq:orientation-isomorphism}
 \det\mathscr D_{\gamma(g)}
 \cong\det T_g\cG^+\otimes(\det\Upsilon_g)^*.
\end{equation}
Under this isomorphism, $D\gamma$ identifies the transverse-zero-locus
orientation of $\chi^{-1}(0)$ with the orientation
$O^{\mathrm{asd}}(\Omega,c_1(L))$ of the actual monopole moduli space.
\end{proposition}

\begin{proof}
The tangent line of the virtual domain is the product of the first four
factors in~\eqref{eq:determinant-splitting}, while
$\Upsilon_g=\Xi_p\oplus K_I$.  This gives~\eqref{eq:orientation-isomorphism}
as a line isomorphism.  To identify it with the finite-dimensional
reduction, differentiate~\eqref{eq:extended-identity} at a zero.  The
derivative of $J_g$ drops out, and
\[
 \dot g\longmapsto(D\gamma_g\dot g,D\chi_g\dot g)
\]
identifies $T_g\cG^+$ with the kernel of the stabilized operator.  It is
homotopic under the model retraction to the identity, with positive Gram
determinant.  The ordering $K_I\oplus\Xi_p$ used in one source formula may be
interchanged with $\Xi_p\oplus K_I$ without sign because both have even real
dimension.  Similarly, the analytic stabilization is $-J_g$, but
$\dim_\RR\Upsilon_g=2r_\Xi+2$ is even.

Proposition~\ref{prop:transversality} gives the exact sequence
\[
 0\longrightarrow T_g\chi^{-1}(0)\longrightarrow T_g\cG^+
 \xrightarrow{D\chi_g}\Upsilon_g\longrightarrow0.
\]
Its determinant is the right side of~\eqref{eq:orientation-isomorphism}
with the same ordered factors.  The claim follows.
\end{proof}

\subsection{The residual circle and the link boundary}
Let $\cZ^+=\chi^{-1}(0)\cap\cG^+$.  On the free circle locus define
\begin{equation}\label{eq:circle-quotient-orientation}
 \lambda_{\cZ}=v_{S^1}\wedge\widetilde\lambda_{\cZ/S^1}.
\end{equation}
The weight-two presentation changes the positive speed of an orbit but not
its orientation.  Let $\vec r$ point outward from the deleted neighborhood
of the ideal stratum.  The boundary of its complement has
\[
 \lambda_{\cZ/S^1}=-\vec r\wedge\lambda_L.
\]
Together with~\eqref{eq:circle-quotient-orientation}, this proves
\begin{equation}\label{eq:link-orientation}
 \lambda_{\cM}=-v_{S^1}\wedge\vec r
 \wedge\widetilde\lambda_L,
\end{equation}
which is the convention in
\cite[Equation~(3.81)]{FeehanLenessLevelOne}.  The sign can also be read in
one complex fiber: moving the positive orbit vector $i\vec r$ past the
outward radial vector contributes the displayed minus sign.

The compactified virtual link has a background-normal face and an
instanton-scale face.  Their common edge receives opposite orientations,
because exchanging the two outward normals contributes a minus sign.  Thus
the faces form an oriented relative cycle.  This statement concerns the
smooth positive-scale zero set and the relative Euler cycle; it does not
reinterpret the lowest stratum as a transverse zero set.

Finally, if $w$ is an integral lift of $w_2(\mathfrak t')$, the orientation
comparison of~\cite{FeehanLenessII} is
\begin{equation}\label{eq:lift-orientation}
 O^{\mathrm{asd}}(\Omega,w)
 =(-1)^{(w-c_1(L))^2/4}
 O^{\mathrm{asd}}(\Omega,c_1(L)).
\end{equation}
The exponent is integral.  The circle and boundary constructions are linear
in the ambient orientation, so the same factor relates the corresponding
link orientations.  This completes the orientation assertion in
Theorem~\ref{thm:neighborhood}.

\section{The level-one pairing}\label{sec:pairing}

We now apply the family neighborhood theorem to the published level-one
pairing.  The neighborhood theorem is used to identify the actual
and virtual links, pull back geometric representatives, carry the universal
bundle through an equivariant homotopy to the splicing map, represent the
smooth zero locus by a relative Euler class, and preserve signed
intersections.  Corollary~\ref{cor:embedding}, Proposition
\ref{prop:neighborhood}, Proposition~\ref{prop:obstruction-section}, and
Proposition~\ref{prop:chart-orientation} supply these inputs.  The dimension
argument of~\cite[Lemmas~3.9 and~5.4]{FeehanLenessLevelOne} puts the counted
intersection in the positive-scale smooth locus and away from the common
edge.  The universal-bundle calculation, Thom and Segre reduction,
charge-one fiber evaluations, and Jacobi identity are retained as published
topological inputs; the contribution here is to discharge their deferred
analytic neighborhood hypothesis.

\subsection{Pullbacks and Euler classes}
Let
\begin{equation}\label{eq:nu-definition}
 \nu=c_1\bigl(\overline\cG^{\mathrm{vir},*}\longrightarrow
 \overline\cG^{\mathrm{vir},*}/S^1\bigr)
\end{equation}
for the effective circle.  The line defining $\mu_c$ on the monopole
configuration quotient has weight $-2$, while the action on the gluing
domain before effectivization has twice the angular parameter.  Hence
\begin{equation}\label{eq:circle-pullback}
 \gamma^*\mu_c=-\nu.
\end{equation}

Away from the incidence diagonal, the pulled-back universal $SO(3)$ bundle
splits as $\RR\oplus\overline L^{\mathrm{vir}}$, where
\begin{equation}\label{eq:associated-line-c1}
 c_1(\overline L^{\mathrm{vir}})=
 2\pi_{\mathfrak s}^*\mu_{\mathfrak s}-\nu
 +\pi_{X,2}^*c_1(L),\qquad c_1(L)=-\beta.
\end{equation}
The equivariant isotopy~\eqref{eq:isotopy} identifies the corrected and
splicing pullbacks.  Insertion of one instanton changes the Pontrjagin class
by the incidence diagonal, so
\begin{align}\label{eq:universal-p1}
 (\gamma\times\id_X)^*p_1(\mathbb F_{\mathfrak t'})
 ={}&\bigl(2\pi_{\mathfrak s}^*\mu_{\mathfrak s}-\nu
 +\pi_{X,2}^*c_1(L)\bigr)^2\notag\\
 &-4(\pi_X\times\id_X)^*\operatorname{PD}[\Delta].
\end{align}
The coefficient $-4$ is fixed by a chain meeting the diagonal once.  Since
$\mu_p(\alpha)=-\tfrac14p_1(\mathbb F_{\mathfrak t'})/\alpha$, slant product
with a point and a surface gives
\begin{align}\label{eq:mu-pullbacks}
 \gamma^*\mu_p(x)
 &=-\tfrac14(2\mu_{\mathfrak s}-\nu)^2
   +\pi_X^*\operatorname{PD}[x],\notag\\
 \gamma^*\mu_p(h)
 &=\tfrac12\langle\beta,h\rangle
   (2\mu_{\mathfrak s}-\nu)+\pi_X^*\operatorname{PD}[h].
\end{align}
These are the formulas of
\cite[Lemma~4.10]{FeehanLenessLevelOne}; the second summand in each line is
the diagonal correction.

The background Euler class is~\eqref{eq:background-euler}.  For the
instanton factor, the square of the framing line loses the central twist:
\[
 \bigl(\SU(2)\times_{\{\pm1\}}\CC\bigr)^{\otimes2}
 \cong\SO(3)\times\CC.
\]
Accounting for the central and residual circle weights gives
\begin{equation}\label{eq:instanton-euler-square}
 c_1\bigl((\Upsilon^i/S^1)^{\otimes2}\bigr)
 =\pi_X^*c_1(\mathfrak t')-\nu.
\end{equation}
Thus, in rational cohomology,
\begin{equation}\label{eq:instanton-euler}
 e(\Upsilon^i/S^1)
 =\tfrac12\bigl(\pi_X^*c_1(\mathfrak t')-\nu\bigr).
\end{equation}
The factor $1/2$ is the square-root factor of the framed associated line; it
has no relation to the derivative-normalization homotopy discussed before
\eqref{eq:determinant-splitting}.

\subsection{Relative Euler reduction}
On the positive-scale virtual link, the section determines a relative Euler
class.  The geometric-representative cocycle is supported away from the
lower-level zero set.  The homotopy
\begin{equation}\label{eq:euler-homotopy}
\chi_s\oplus\chi_i\longmapsto\chi_s\oplus t\chi_i,
\qquad 0\leq t\leq1,
\end{equation}
does not vanish off $\chi_s^{-1}(0)$.  At $t=0$ the class factors as
\begin{equation}\label{eq:euler-factorization}
 e(\Upsilon/S^1,\chi_s)
 =e(\overline\Upsilon^s/S^1,\chi_s)
  \smile e(\Upsilon^i/S^1).
\end{equation}
This is precisely where one needs continuity of the background section and
the cone topology, but not transversality of the lowest stratum.  The signed
actual-link count is therefore
\begin{equation}\label{eq:relative-pairing}
 \left\langle
 \overline\mu_p(z)\smile\overline\mu_c^{\eta}
 \smile\overline e_s\smile\overline e_i,
 [\overline L^{\mathrm{vir}}]\right\rangle.
\end{equation}

Write
\[
 H=\pi_X^*\operatorname{PD}[h],\qquad
 X_0=\pi_X^*\operatorname{PD}[x],\qquad
 c=\pi_X^*c_1(\mathfrak t').
\]
Substitution of~\eqref{eq:circle-pullback},
\eqref{eq:mu-pullbacks},~\eqref{eq:background-euler}, and
\eqref{eq:instanton-euler} into~\eqref{eq:relative-pairing} gives
\begin{align}\label{eq:pairing-integrand}
 \mathcal P=\Bigl\langle&
 \left[\tfrac12B_h(2\mu_{\mathfrak s}-\nu)+H\right]^n\notag\\
 &\smile\left[-\tfrac14(2\mu_{\mathfrak s}-\nu)^2+X_0\right]^m
 (-\nu)^{\eta+r_\Xi}\notag\\
 &\smile\tfrac12(c-\nu),[\overline L^{\mathrm{vir}}]
 \Bigr\rangle.
\end{align}
The dimension relation is
\begin{equation}\label{eq:pairing-dimension}
 \delta+\eta=d+r_N-r_\Xi+2.
\end{equation}

The equivariant Thom division of
\cite[Propositions~5.20--5.21]{FeehanLenessLevelOne} expresses this pairing
over
\[
 M_{\mathfrak s}\times
 \partial\overline{\Gl}_{\mathfrak t}(\delta)/S^1
\]
as a finite sum involving the Segre classes of $N$.  The normal-bundle
calculation~\cite[Lemma~2.4]{FeehanLenessLevelOne} gives
\begin{equation}\label{eq:normal-segre}
 s_j(N)=S_j\mu_{\mathfrak s}^{j},
\end{equation}
where $S_j$ is the binomial coefficient in
\cite[Equation~(6.14)]{FeehanLenessLevelOne}.  On this product base,
\cite[Lemma~5.23]{FeehanLenessLevelOne} gives
\begin{equation}\label{eq:circle-splitting}
 \nu=\nu_{\mathfrak t}+2\mu_{\mathfrak s}.
\end{equation}
Every nonzero term then has base factor
\begin{equation}\label{eq:sw-base-pairing}
 \left\langle\mu_{\mathfrak s}^{d},[M_{\mathfrak s}]\right\rangle
 =\SW_X(\mathfrak s).
\end{equation}
In particular, positive-dimensional base classes have not been discarded;
the Segre calculation combines all of them into the single pairing in
\eqref{eq:sw-base-pairing}.

Put
\begin{equation}\label{eq:jacobi-parameters}
 A_J=\eta-d+1,\qquad
 B_J=\tfrac12(2\delta-d_a(\mathfrak t')-d_{\mathfrak s})
      -\tfrac14\bigl(\chi(X)+\sigma(X)\bigr),
\end{equation}
and set
\begin{equation}\label{eq:jacobi-values}
 J_d=P_d^{A_J,B_J}(0),\qquad
 J_d'=P_d^{A_J-1,B_J+1}(0),
\end{equation}
using~\eqref{eq:intro-jacobi-normalization}.  The finite Segre sum is
$2^dJ_d$, except in the term containing $Hc$, where it is $2^dJ_d'$;
this is the Jacobi identity of
\cite[Lemma~6.2]{FeehanLenessLevelOne}.

\begin{lemma}[published charge-one link integrals]
\label{lem:instanton-link-integrals}
Let $\nu_t$ denote the effective circle class on the charge-one instanton
link.  With the orientation in~\eqref{eq:link-orientation},
\begin{align}\label{eq:instanton-link-integrals}
 \langle\nu_tX_0,[\partial\overline\Gl/S^1]\rangle&=2,
 &\langle\nu_t^2H,[\partial\overline\Gl/S^1]\rangle&=-4B_h,
 \notag\\
 \langle\nu_tHc,[\partial\overline\Gl/S^1]\rangle
 &=2\langle c_1(\mathfrak t'),h\rangle,
 &\langle\nu_t^3,[\partial\overline\Gl/S^1]\rangle
 &=6\beta^2+2c_1^2(X),\notag\\
 \langle\nu_t^2c,[\partial\overline\Gl/S^1]\rangle
 &=-4\beta\cdot c_1(\mathfrak t').&&
\end{align}
\end{lemma}

\begin{proof}
These numerical evaluations are cohomological inputs from
\cite[Lemma~5.24 and Equation~(5.77)]{FeehanLenessLevelOne}; they are not
reproved here.  We recall enough of that calculation to fix their
normalization.  On a finite cover of $X$, the framed charge-one gluing bundle
is covered with degree $-2$ by the boundary of a rank-two Hermitian bundle
$F$.  The induced map of projectivizations has degree $-1$, while $\nu_t$
pulls back to twice the hyperplane class.  The projective-bundle pushforward
\[
 \pi_*(h^{1+j})=s_j(F)
\]
therefore gives
$\langle\nu_tX_0\rangle=2$ and, from the published Segre classes of $F$,
the second and fourth values in~\eqref{eq:instanton-link-integrals}.  Taking
$h=\operatorname{PD}(c)$ in the second value gives the fifth.  Finally,
Poincare duality in Equation~(5.77) gives
\[
 \langle\nu_tHc\rangle
 =\langle\nu_tX_0\rangle\langle c_1(\mathfrak t'),h\rangle,
\]
which is the third value.  Since $c_1(L)=-\beta$, these conventions agree
with~\eqref{eq:circle-pullback} and~\eqref{eq:link-orientation}.
\end{proof}

Combining the base reduction with
Lemma~\ref{lem:instanton-link-integrals} gives the following polynomial form,
which does not divide by a Jacobi value and therefore remains defined when
$J_d=0$.

\begin{corollary}[completion of the family level-one pairing]\label{cor:pairing}
Assume the four-manifold, level-one, regularity, and nonzero-spinor hypotheses
of Theorem~\ref{thm:neighborhood}.  Suppose in addition that
$w_2(\mathfrak t')$ is good in the sense of
\cite[Definition~2.3]{FeehanLenessLevelOne}.  Let $\delta,m,\eta$ be
nonnegative integers satisfying
\[
 0\leq m\leq\lfloor\delta/2\rfloor,
 \qquad
 2(\delta+\eta)=\dim(\cM_{\mathfrak t'}/S^1)-1,
\]
and use the notation in
\eqref{eq:intro-pairing-notation} and
\eqref{eq:jacobi-parameters}--\eqref{eq:jacobi-values}.  Then
\begin{align*}
&\#\bigl(\overline V(h^{\delta-2m}x^m)\cap\overline W^\eta
\cap\overline L_{\mathfrak t',\mathfrak s}\bigr)\\
&\quad=(-1)^{m+1+d}2^{d-\delta}
\left\langle\mu_{\mathfrak s}^{d},[M_{\mathfrak s}]\right\rangle\\
&\qquad\quad\times
\left[J_d\left(a_0B_h^n+a_1B_h^{n-2}Q_X(h,h)\right)
+2nJ_d'B_h^{n-1}\langle c_1(\mathfrak t'),h\rangle\right],
\end{align*}
where terms with negative powers of $B_h$ are omitted and
\begin{align*}
 a_0&=3\beta^2+c_1^2(X)+2\beta\cdot c_1(\mathfrak t')
 +4n-4m,\qquad
 a_1=4\binom n2.
\end{align*}
\end{corollary}

\begin{proof}
Equation~\eqref{eq:pairing-integrand}, equivariant Thom division,
\eqref{eq:normal-segre}, and~\eqref{eq:circle-splitting} reduce the base
degree to~\eqref{eq:sw-base-pairing}.  The Jacobi identity gives $2^dJ_d$
and $2^dJ_d'$ in the two indicated groups of terms.  The five charge-one
evaluations in Lemma~\ref{lem:instanton-link-integrals} give respectively
the $-4m$, $4n$, $2\beta\cdot c_1(\mathfrak t')$, $c_1^2(X)$, and
$Q_X(h,h)$ contributions.  Collecting them gives $a_0$ and $a_1$.
The sign and power of two are
$(-1)^{m+1+d}2^{d-\delta}$, as in
\cite[Equations~(6.15)--(6.18)]{FeehanLenessLevelOne}; Proposition
\ref{prop:chart-orientation} introduces no further sign.  This proves the
formula.
\end{proof}

\begin{corollary}[zero-dimensional specialization]
\label{cor:zero-dimensional-pairing}
If $d_{\mathfrak s}=0$, then $J_0=J_0'=1$ and
\begin{align*}
&\#\bigl(\overline V(h^{\delta-2m}x^m)\cap\overline W^\eta
\cap\overline L_{\mathfrak t',\mathfrak s}\bigr)\\
&\quad=(-1)^{m+1}2^{-\delta}\SW_X(\mathfrak s)
\left(a_0B_h^n+b_0B_h^{n-1}\langle c_1(\mathfrak t'),h\rangle
+a_1B_h^{n-2}Q_X(h,h)\right),
\end{align*}
where $b_0=2n$ and $a_0,a_1$ are as in Corollary~\ref{cor:pairing}.
\end{corollary}

\subsection{Blow-up and the degree consequence}
The only blow-up specialization needed below has no hidden change in
$a_0$.  On $\widetilde X=X\#\overline{\mathbb{CP}}^{\,2}$, let $e$ be the
exceptional class and
\[
 c_1(\mathfrak s^\pm)=c_1(\mathfrak s)\pm\operatorname{PD}[e],\qquad
 \SW_{\widetilde X}(\mathfrak s^\pm)=\SW_X(\mathfrak s).
\]
For $\widetilde{\mathfrak t}'$ one has
$c_1(\widetilde{\mathfrak t}')=c_1(\mathfrak t')$ and
$p_1(\widetilde{\mathfrak t}')=p_1(\mathfrak t')-1$.  The exceptional
evaluations are $\mp1$, while~\eqref{eq:lift-orientation} gives exponents
$o_{\widetilde{\mathfrak t}'}(w+\operatorname{PD}[e],\mathfrak s^+)
=o_{\mathfrak t'}(w,\mathfrak s)-1$ and
$o_{\widetilde{\mathfrak t}'}(w+\operatorname{PD}[e],\mathfrak s^-)
=o_{\mathfrak t'}(w,\mathfrak s)$.  The two exceptional contributions
therefore add after orientation weighting.  Their factor two converts
$2^{-(\delta+1)}$ to $2^{-\delta}$.

Algebraically,
\begin{align*}
&3(\beta^2-1)+(c_1^2(X)-1)+2\beta\cdot c_1(\mathfrak t')
 +4(\delta+1-2m)-4m\\
&\qquad=3\beta^2+c_1^2(X)+2\beta\cdot c_1(\mathfrak t')
 +4(\delta-2m)-4m.
\end{align*}
Together with $Q_{\widetilde X}(e,h)=0$ and $e^2=-1$, this is the $k=0$
case of~\cite[Proposition~6.5]{FeehanLenessLevelOne}.

\begin{corollary}[degree consequence]\label{cor:degree}
Suppose in addition that $X$ is abundant, effective, and of
Seiberg--Witten simple type, and choose the characteristic data as in
\cite[Theorem~1.1]{FeehanLenessLevelOne}.  Then
\begin{align}\label{eq:series-congruence}
 \cD_X^w(h)&\equiv0\equiv\SW_X^w(h)
 &&\pmod{h^{c(X)-2}},\notag\\
 \cD_X^w(h)&\equiv2^{2-c(X)}e^{Q_X(h,h)/2}\SW_X^w(h)
 &&\pmod{h^{c(X)+2}}.
\end{align}
In the last two relevant degrees,
\begin{align}\label{eq:terminal-degrees}
 D_X^w(h^{c(X)-2}x)
 &=2^{3-c(X)}\sum_{\mathfrak s}
 (-1)^{(w^2+w\cdot c_1(\mathfrak s))/2}\SW_X(\mathfrak s)
 \langle c_1(\mathfrak s),h\rangle^{c(X)-2},\notag\\
 D_X^w(h^{c(X)})
 &=2^{2-c(X)}\sum_{\mathfrak s}
 (-1)^{(w^2+w\cdot c_1(\mathfrak s))/2}\SW_X(\mathfrak s)\notag\\
 &\quad\times\left(\langle c_1(\mathfrak s),h\rangle^{c(X)}
 +\binom{c(X)}2\langle c_1(\mathfrak s),h\rangle^{c(X)-2}
 Q_X(h,h)\right).
\end{align}
\end{corollary}

\begin{proof}
For abundant, effective manifolds of Seiberg--Witten simple type, choose the
class $\Lambda$ and lift $w$ as in
\cite[Theorem~1.1]{FeehanLenessLevelOne}.  Then
$r(\Lambda)=c(X)-4$ and $i(\Lambda)=c(X)+4$.  The published top-level
pairing~\cite{FeehanLenessII}, Corollary~\ref{cor:pairing}, the Donaldson and
Seiberg--Witten blow-up formulas, and the superconformal-simple-type
vanishing used in~\cite[Section~6]{FeehanLenessLevelOne} give
\eqref{eq:series-congruence}.  Extracting the coefficients of the two final
degrees gives~\eqref{eq:terminal-degrees}.
\end{proof}

As a normalization test, for a K3 surface one has $c(X)=2$ and only the zero
basic class, with Seiberg--Witten invariant one in the standard homology
orientation.  Formula~\eqref{eq:terminal-degrees} becomes
\begin{equation}\label{eq:k3-test}
 D_{K3}(x)=2,\qquad D_{K3}(h^2)=Q_{K3}(h,h),
\end{equation}
the constant and quadratic coefficients of $e^{Q_{K3}(h,h)/2}$ in the
convention with insertion $1+x/2$.  Thus the global test detects neither an
additional sign nor an additional factor two.

\section{Scope and further problems}\label{sec:discussion}

Theorem~\ref{thm:neighborhood} supplies the analytic and orientation input
that was deferred in the one-bubble calculation of
\cite{FeehanLenessLevelOne}: the corrected chart is continuous and embedded,
its zero set is an actual Uhlenbeck neighborhood, and its determinant and
boundary orientations agree with those used in the pairing.  The proof also
shows why the compactified instanton obstruction must be treated as a cone
pseudo-bundle and why transversality stops at the middle stratum.

The compact background is permitted to have arbitrary dimension, but its
regularity and nonzero-spinor hypotheses are essential here.  The proof does
not treat a zero-spinor pair quotient or a nonregular base.  It avoids the
varying-rank cutoff problem by global finite-dimensional stabilization;
spectral bundles are used only on positive-scale charts where their rank is
locally constant.

There is exactly one charge-one bubble.  Consequently there are no
collision diagonals, no competing partitions of a symmetric product, and no
compatibility maps between gluing charts.  The topological overlap machinery
of~\cite{FeehanLenessCobordism} addresses those phenomena, but its analytic
neighborhood identification remains a separate issue.  The present theorem
does not imply the corresponding result at level two or higher.

Corollary~\ref{cor:degree} has an equally precise logical range.  It removes
the deferred neighborhood theorem from the level-one contributions.  It
continues to use the published top-level pairing, the
Donaldson and Seiberg--Witten blow-up formulas, the
superconformal-simple-type vanishing input, and the assumptions of abundance,
effectiveness, and Seiberg--Witten simple type.  In particular, it is neither
a proof of the full Witten conjecture nor a removal of the effectiveness
hypothesis.

\begin{thebibliography}{FL18}

\bibitem[Don87]{DonaldsonOrientation} S.~K. Donaldson,
\emph{The orientation of Yang--Mills moduli spaces and $4$-manifold
topology}, J.\ Differential Geom.\ \textbf{26} (1987), no.~3, 397--428.

\bibitem[DK90]{DK90} S.~K. Donaldson and P.~B. Kronheimer,
\emph{The geometry of four-manifolds}, Oxford Mathematical Monographs,
Clarendon Press, Oxford, 1990.

\bibitem[Fee01]{FeehanCritical} P.~M.~N. Feehan,
\emph{Critical-exponent Sobolev norms and the slice theorem for the quotient
space of connections}, Pacific J.\ Math.\ \textbf{200} (2001), no.~1,
71--118; arXiv:dg-ga/9711004.

\bibitem[FL98]{FeehanLenessI} P.~M.~N. Feehan and T.~G. Leness,
\emph{$PU(2)$ monopoles. I: Regularity, Uhlenbeck compactness, and
transversality}, J.\ Differential Geom.\ \textbf{49} (1998), no.~2,
265--410; arXiv:dg-ga/9710032.

\bibitem[FL01a]{FeehanLenessTopLevel} P.~M.~N. Feehan and T.~G. Leness,
\emph{$PU(2)$ monopoles and links of top-level Seiberg--Witten moduli
spaces}, J.\ Reine Angew.\ Math.\ \textbf{538} (2001), 57--133;
arXiv:math/0007190.

\bibitem[FL01b]{FeehanLenessII} P.~M.~N. Feehan and T.~G. Leness,
\emph{$PU(2)$ monopoles. II: Top-level Seiberg--Witten moduli spaces and
Witten's conjecture in low degrees}, J.\ Reine Angew.\ Math.\ \textbf{538}
(2001), 135--212; arXiv:dg-ga/9712005.

\bibitem[FL99]{FeehanLenessIII} P.~M.~N. Feehan and T.~G. Leness,
\emph{$PU(2)$ monopoles. III: Existence of gluing and obstruction maps},
preprint, 1999; arXiv:math/9907107.

\bibitem[FL02]{FeehanLenessLevelOne} P.~M.~N. Feehan and T.~G. Leness,
\emph{$SO(3)$ monopoles, level-one Seiberg--Witten moduli spaces, and
Witten's conjecture in low degrees}, Topology Appl.\ \textbf{124} (2002),
no.~2, 221--326; arXiv:math/0106238.

\bibitem[FL18]{FeehanLenessCobordism} P.~M.~N. Feehan and T.~G. Leness,
\emph{An $SO(3)$-monopole cobordism formula relating Donaldson and
Seiberg--Witten invariants}, Mem.\ Amer.\ Math.\ Soc.\ \textbf{256} (2018),
no.~1226; arXiv:math/0203047.

\bibitem[Uhl82]{Uhlenbeck82} K.~K. Uhlenbeck,
\emph{Connections with $L^p$ bounds on curvature}, Comm.\ Math.\ Phys.\
\textbf{83} (1982), no.~1, 31--42.

\end{thebibliography}

\bigskip
\noindent{\sc Bernd Johannes Wuebben, New York, NY,}
\texttt{wuebben@gmail.com}

\end{document}
